\documentclass[../main.tex]{subfiles}
\begin{document}
本章主要参考\cite{dlmf,hunter}。

\section{函数、极限、连续}
\subsection{基本初等函数}\index{基本初等函数|hyperpage}
幂函数 $y=x^\alpha$（实数幂通常取 $x>0$）；指数函数 $y=a^x$、对数函数 $y=\log_a x$（$a>0,a\ne1$，对数要求 $x>0$）。
\[
a^{x+y}=a^xa^y,\quad \log_a(xy)=\log_a x+\log_a y,\quad
\log_a x=\frac{\ln x}{\ln a}.
\]
三角函数及其反函数见初等数学章。双曲函数满足
\[
\sinh x=\frac{e^x-e^{-x}}2,\quad \cosh x=\frac{e^x+e^{-x}}2,\quad
\tanh x=\frac{\sinh x}{\cosh x},\quad \cosh^2x-\sinh^2x=1.
\]
\subsection{函数的性质}\index{函数的性质|hyperpage}
对定义域内任意 $x_1<x_2$，$f(x_1)\le f(x_2)$ 表示单调不减；严格小于表示严格递增。
\begin{align*}
\text{有界性：}&\quad \exists M>0,\ \forall x\in D,\ |f(x)|\le M,\\
\text{奇偶性：}&\quad f(-x)=-f(x)\ \text{（奇）},\quad f(-x)=f(x)\ \text{（偶）},\\
\text{周期性：}&\quad \exists T>0,\ f(x+T)=f(x).
\end{align*}
奇偶性要求定义域关于原点对称；周期性要求定义域在平移 $T$ 下保持不变。
\subsection{数列极限的定义及其性质}\index{数列极限的定义及其性质|hyperpage}
\[
\lim_{n\to\infty}a_n=A\iff
\forall\varepsilon>0,\ \exists N,\ \forall n>N,\ |a_n-A|<\varepsilon.
\]
收敛数列有界，极限唯一，任意子列收敛于同一极限。若 $a_n\to A,b_n\to B$，则
\[
a_n\pm b_n\to A\pm B,\quad a_nb_n\to AB,\quad
\frac{a_n}{b_n}\to\frac AB\ (B\ne0).
\]
夹逼定理：$a_n\le b_n\le c_n$ 且 $a_n,c_n\to L$，则 $b_n\to L$。
单调有界数列必收敛，且
\[
\left(1+\frac1n\right)^n\to e,\qquad q^n\to0\quad(|q|<1).
\]
\subsection{Cauchy准则、子列与上下极限}\index{Cauchy准则、子列与上下极限|hyperpage}
实数列收敛当且仅当它是 Cauchy 列：
\[
\forall\varepsilon>0,\ \exists N,\quad
m,n\ge N\Longrightarrow |a_m-a_n|<\varepsilon.
\]
Bolzano--Weierstrass 定理：有界实数列必有收敛子列。对有界实数列，若所有收敛子列的极限均为 $A$，则原数列收敛于 $A$。
在扩展实数意义下定义
\[
\overline a=\limsup_{n\to\infty}a_n
=\lim_{n\to\infty}\sup_{k\ge n}a_k,\qquad
\underline a=\liminf_{n\to\infty}a_n
=\lim_{n\to\infty}\inf_{k\ge n}a_k.
\]
其中尾上确界单调不增，尾下确界单调不减，且 $\underline a\le\overline a$。对有界实数列，二者分别为子列极限的最小值与最大值，并且
\[
a_n\to A\in\mathbb R\iff\underline a=\overline a=A.
\]
\subsection{Stolz--Ces\`aro定理与平均极限}\index{Stolz-Cesàro定理与平均极限|hyperpage}
若 $b_n$ 严格递增且 $b_n\to+\infty$，则对 $L\in\mathbb R\cup\{\pm\infty\}$，
\[
\frac{a_{n+1}-a_n}{b_{n+1}-b_n}\to L
\Longrightarrow\frac{a_n}{b_n}\to L.
\]
零比零形式：若 $a_n\to0$，$b_n>0$ 严格递减至零，则上述差商收敛于 $L$ 时，仍有 $a_n/b_n\to L$。
作为平均极限公式，对有限实数 $A$，
\[
a_n\to A\Longrightarrow
\frac1n\sum_{k=1}^n a_k\to A;
\qquad
a_n>0,\ a_n\to A>0\Longrightarrow
\left(\prod_{k=1}^na_k\right)^{1/n}\to A.
\]
\subsection{函数极限的概念及其关系}\index{函数极限的概念及其关系|hyperpage}
\[
\lim_{x\to x_0}f(x)=A\iff\forall\varepsilon>0,\ \exists\delta>0:
0<|x-x_0|<\delta\Longrightarrow |f(x)-A|<\varepsilon.
\]
双侧极限存在当且仅当左右极限存在且相等。Heine 定理：上述极限等价于对任意定义域中 $x_n\ne x_0$、$x_n\to x_0$ 均有 $f(x_n)\to A$。
\[
\lim_{x\to0}\frac{\sin x}{x}=1,\quad
\lim_{x\to0}\frac{e^x-1}{x}=1,\quad
\lim_{x\to0}\frac{\ln(1+x)}x=1,\quad
\lim_{x\to0}(1+x)^{1/x}=e.
\]
L'Hospital 法则：在 $0/0$ 或 $\infty/\infty$ 型中，若去心单侧邻域内 $f,g$ 可导、$g'\ne0$，且 $f'/g'$ 有有限或无穷极限，则 $\lim f/g=\lim f'/g'$。
\subsection{无穷小和无穷大的概念及其关系}\index{无穷小和无穷大的概念及其关系|hyperpage}
在指定极限过程中，$\alpha\to0$ 称为无穷小，$|\beta|\to\infty$ 称为无穷大；非零时二者互为倒数。
\[
\alpha=o(\beta)\iff\frac\alpha\beta\to0,\quad
\alpha\sim\beta\iff\frac\alpha\beta\to1,\quad
\alpha=O(\beta)\iff\left|\frac\alpha\beta\right|\text{在该过程中有界}.
\]
当 $x\to0$，有
\[
\sin x\sim\tan x\sim\arcsin x\sim\arctan x\sim x,\qquad
1-\cos x\sim\frac{x^2}2,
\]
\[
e^x-1\sim\ln(1+x)\sim x,\qquad (1+x)^\alpha-1\sim\alpha x\quad(\alpha\ne0).
\]
等价无穷小可用于乘除中的替换，和差中应先检查抵消项。
\subsection{连续与间断}\index{连续与间断|hyperpage}
\[
f\text{ 在 }x_0\text{ 连续}\iff\lim_{x\to x_0}f(x)=f(x_0).
\]
左右极限有限且相等但不等于函数值（或函数值未定义）为可去间断；左右极限有限但不等为跳跃间断；至少一侧无有限极限为第二类间断。
闭区间上的连续函数有界、取得最大最小值，满足介值定理并一致连续：
\[
\forall\varepsilon>0,\ \exists\delta>0,\ \forall x,y\in[a,b],\quad
|x-y|<\delta\Longrightarrow |f(x)-f(y)|<\varepsilon.
\]
\subsection{连续模、Lipschitz与H\"older条件}\index{连续模、Lipschitz与Hölder条件|hyperpage}
设 $E\subset\mathbb R$ 非空，$f:E\to\mathbb R$。连续模定义为
\[
\omega_f(\delta)=\sup_{\substack{x,y\in E\\|x-y|\le\delta}}
|f(x)-f(y)|\quad(\delta\ge0),
\]
上确界允许为 $+\infty$。$\omega_f$ 单调不减，且
\[
f\text{ 在 }E\text{ 上一致连续}
\iff\lim_{\delta\downarrow0}\omega_f(\delta)=0.
\]
若存在 $C\ge0$、$0<\alpha\le1$，使
\[
|f(x)-f(y)|\le C|x-y|^\alpha\quad(x,y\in E),
\]
则称 $f$ 满足 $\alpha$ 阶 H\"older 条件；$\alpha=1$ 时为 Lipschitz 条件。此时 $\omega_f(\delta)\le C\delta^\alpha$，因而一致连续。
若 $E$ 的直径 $d\in(0,\infty)$，且 $0<\alpha<\beta\le1$，则
\[
|f(x)-f(y)|\le C|x-y|^\beta
\Longrightarrow |f(x)-f(y)|\le Cd^{\beta-\alpha}|x-y|^\alpha.
\]
对区间 $I$ 上的可导函数，中值定理给出
\[
\sup_{x\in I}|f'(x)|\le M
\Longrightarrow |f(x)-f(y)|\le M|x-y|\quad(x,y\in I).
\]
\section{一元函数微分学}
\[
f'(x)=\lim_{h\to0}\frac{f(x+h)-f(x)}h,\qquad \mathrm df=f'(x)\,\mathrm dx.
\]
\begin{align*}
(x^\alpha)'&=\alpha x^{\alpha-1}, &(e^x)'&=e^x, &(a^x)'&=a^x\ln a,\\
(\ln|x|)'&=1/x, &(\sin x)'&=\cos x, &(\cos x)'&=-\sin x,\\
(\tan x)'&=\sec^2x, &(\cot x)'&=-\csc^2x,\\
(\arcsin x)'&=\frac1{\sqrt{1-x^2}}, &(\arccos x)'&=-\frac1{\sqrt{1-x^2}},\\
(\arctan x)'&=\frac1{1+x^2}, &(\operatorname{arccot}x)'&=-\frac1{1+x^2}.
\end{align*}
以上公式在相应函数可导的定义域内使用。
\[
(uv)'=u'v+uv',\quad (u/v)'=\frac{u'v-uv'}{v^2},\quad
(f\circ g)'=(f'\circ g)g',\quad (f^{-1})'(y)=\frac1{f'(x)}.
\]
商要求 $v\ne0$，反函数公式要求局部可逆且 $f'(x)\ne0$。
\subsection{常用导数补充与高阶求导}\index{常用导数补充与高阶求导|hyperpage}
以下公式在各函数可导区间内成立：
\begin{align*}
(\log_a x)'&=\frac1{x\ln a}\quad(x>0,\ a>0,\ a\ne1),\\
(\sec x)'&=\sec x\tan x, &(\csc x)'&=-\csc x\cot x,\\
(\sinh x)'&=\cosh x, &(\cosh x)'&=\sinh x,\\
(\tanh x)'&=\frac1{\cosh^2x},
 &(\operatorname{arsinh}x)'&=\frac1{\sqrt{1+x^2}}.
\end{align*}
若 $u,v$ 有 $n$ 阶导数，Leibniz 公式为
\[
(uv)^{(n)}=\sum_{k=0}^n\binom nk u^{(k)}v^{(n-k)}.
\]
\begin{align*}
(e^{ax})^{(n)}&=a^ne^{ax},\\
(\sin ax)^{(n)}&=a^n\sin(ax+n\pi/2),\\
(\cos ax)^{(n)}&=a^n\cos(ax+n\pi/2),\\
(\ln|x|)^{(n)}&=(-1)^{n-1}(n-1)!/x^n\quad(n\ge1,\ x\ne0).
\end{align*}
对 $y=u(x)^{v(x)}$、$u>0$，对数求导给出 $y'=y[v'\ln u+vu'/u]$。
\subsection{高阶复合求导与Fa\`a di Bruno公式}\index{高阶复合求导与Faà di Bruno公式|hyperpage}
设 $f,g$ 在所讨论点的相应邻域内为 $C^n$，$n\ge1$。令
\[
\mathcal K_n=\left\{(k_1,\ldots,k_n)\in\mathbb N_0^n:
\sum_{j=1}^njk_j=n\right\},\qquad K=\sum_{j=1}^nk_j.
\]
高阶链式法则为
\[
(f\circ g)^{(n)}(x)=
\sum_{\bm k\in\mathcal K_n}
\frac{n!}{k_1!\cdots k_n!}\,f^{(K)}(g(x))
\prod_{j=1}^n\left(\frac{g^{(j)}(x)}{j!}\right)^{k_j}.
\]
常用低阶形式为（右端 $f$ 的各阶导数均在 $g(x)$ 处取值）
\begin{align*}
(f\circ g)''&=f''(g)(g')^2+f'(g)g'',\\
(f\circ g)'''&=f'''(g)(g')^3+3f''(g)g'g''+f'(g)g''',\\
(f\circ g)^{(4)}&=f^{(4)}(g)(g')^4+6f'''(g)(g')^2g''\\
&\quad+3f''(g)(g'')^2+4f''(g)g'g'''+f'(g)g^{(4)}.
\end{align*}
\subsection{反函数的二阶与三阶导数}\index{反函数的二阶与三阶导数|hyperpage}
设 $f\in C^3(I)$，$f'$ 在区间 $I$ 上处处非零，$g=f^{-1}$。对 $y=f(x)$，
\begin{align*}
g'(y)&=\frac1{f'(x)},\\
g''(y)&=-\frac{f''(x)}{[f'(x)]^3},\\
g'''(y)&=\frac{3[f''(x)]^2-f'(x)f'''(x)}{[f'(x)]^5}.
\end{align*}
若 $f\in C^n$，更高阶导数可通过递推算子计算：
\[
g^{(n)}(y)=
\left(\frac1{f'(x)}\frac{\mathrm d}{\mathrm dx}\right)^{n-1}
\frac1{f'(x)},\qquad x=g(y).
\]
\subsection{Taylor公式}\index{Taylor公式|hyperpage}
若 $f$ 在 $a$ 附近有 $n+1$ 阶连续导数，则存在介于 $a,x$ 之间的 $\xi$，使
\[
f(x)=\sum_{k=0}^n\frac{f^{(k)}(a)}{k!}(x-a)^k+
\frac{f^{(n+1)}(\xi)}{(n+1)!}(x-a)^{n+1}.
\]
若 $f$ 在 $a$ 处 $n$ 阶可导，则 Peano 余项形式为
\[
f(x)=\sum_{k=0}^n\frac{f^{(k)}(a)}{k!}(x-a)^k+o((x-a)^n).
\]
\subsection{积分型余项与Taylor误差估计}\index{积分型余项与Taylor误差估计|hyperpage}
设 $n\ge0$，$f\in C^{n+1}(I)$，$a,x\in I$，记
$T_n(x)=\sum_{k=0}^nf^{(k)}(a)(x-a)^k/k!$。积分型余项为
\[
f(x)=T_n(x)+R_n(x),\qquad
R_n(x)=\frac1{n!}\int_a^x(x-t)^nf^{(n+1)}(t)\,\mathrm dt.
\]
若连接 $a,x$ 的区间上 $|f^{(n+1)}(t)|\le M$，则
\[
|R_n(x)|\le\frac{M|x-a|^{n+1}}{(n+1)!}.
\]
若该导数界在 $[a-r,a+r]$ 上成立，则
\[
\sup_{|x-a|\le r}|f(x)-T_n(x)|\le\frac{Mr^{n+1}}{(n+1)!}.
\]
Taylor 级数等于原函数还要求 $R_n(x)\to0$；仅有各阶导数存在不足以保证这一点。
\subsection{高阶导数的极值判别}\index{高阶导数的极值判别|hyperpage}
设 $a$ 为内点，$f\in C^m$ 于 $a$ 的邻域，$m\ge2$，并且
\[
f'(a)=\cdots=f^{(m-1)}(a)=0,\qquad f^{(m)}(a)\ne0.
\]
\par\pagebreak[0]\begin{samepage}
则
\[
f(a+h)-f(a)=\frac{f^{(m)}(a)}{m!}h^m+o(|h|^m).
\]
$m$ 为偶数且 $f^{(m)}(a)>0$ 时为严格极小点，$f^{(m)}(a)<0$ 时为严格极大点；$m$ 为奇数时不是极值点。若所有已知阶导数均为零，此判别不能得出结论。
\par\end{samepage}\pagebreak[0]
\subsection{中值定理}\index{中值定理|hyperpage}
设 $f$ 在 $[a,b]$ 连续、在 $(a,b)$ 可导。Lagrange 定理为
\[
f(b)-f(a)=f'(\xi)(b-a),\qquad \xi\in(a,b).
\]
若 $f(a)=f(b)$，则有 $f'(\xi)=0$（Rolle 定理）。若 $g$ 也满足同样条件，则 Cauchy 定理为
\[
[f(b)-f(a)]g'(\xi)=[g(b)-g(a)]f'(\xi).
\]
\subsection{平面曲线的切线与法线}\index{平面曲线的切线与法线|hyperpage}
在 $(x_0,y_0)$ 处斜率 $k=f'(x_0)$ 有限时，
\[
\text{切线： }y-y_0=k(x-x_0),\qquad
\text{法线： }x-x_0+k(y-y_0)=0.
\]
参数曲线 $x=x(t),y=y(t)$ 在 $x'(t)\ne0$ 时
\[
\frac{\mathrm dy}{\mathrm dx}=\frac{y'}{x'},\qquad
\frac{\mathrm d^2y}{\mathrm dx^2}=\frac{x'y''-y'x''}{(x')^3}.
\]
正则隐曲线 $F(x,y)=0$ 的法向量为 $(F_x,F_y)$，切线为 $F_x(x-x_0)+F_y(y-y_0)=0$（偏导在该点取值）。
\subsection{曲线的凹凸性与拐点}\index{曲线的凹凸性与拐点|hyperpage}
采用 $f''>0$ 为凸（图像在切线上方）、$f''<0$ 为凹的约定。
\[
f''(x)>0\Longrightarrow f'\text{ 递增},\qquad
f''(x)<0\Longrightarrow f'\text{ 递减}.
\]
连续曲线在 $x_0$ 两侧凹凸性改变时，$(x_0,f(x_0))$ 为拐点；若 $f''(x_0)$ 存在，则拐点必有 $f''(x_0)=0$，反之不一定成立。
\subsection{渐近线}\index{渐近线|hyperpage}
\begin{align*}
\lim_{x\to a^\pm}|f(x)|=\infty&\Longrightarrow x=a\text{ 为垂直渐近线},\\
\lim_{x\to\pm\infty}f(x)=b&\Longrightarrow y=b\text{ 为水平渐近线},\\
k=\lim_{x\to\pm\infty}\frac{f(x)}x,\quad
b=\lim_{x\to\pm\infty}[f(x)-kx]&\Longrightarrow y=kx+b\text{ 为渐近线}.
\end{align*}
后一式要求两个极限有限；$k\ne0$ 时为斜渐近线，正负无穷分别考察。
\subsection{曲率}\index{曲率|hyperpage}
\[
\kappa=\frac{|f''(x)|}{[1+f'(x)^2]^{3/2}},\qquad
\kappa=\frac{|x'y''-y'x''|}{(x'^2+y'^2)^{3/2}},\qquad \rho=\frac1\kappa.
\]
参数形式要求 $x'^2+y'^2>0$；曲率半径公式要求 $\kappa>0$。
\section{一元函数积分学}
\subsection{不定积分的基本概念}\index{不定积分的基本概念|hyperpage}
在区间内 $F'=f$ 时，$\int f(x)\,\mathrm dx=F(x)+C$，其中 $C$ 为任意常数。
\begin{align*}
\int x^\alpha\,\mathrm dx&=\frac{x^{\alpha+1}}{\alpha+1}+C\quad(\alpha\ne-1),\\
\int\frac{\mathrm dx}{x}&=\ln|x|+C, & \int e^x\,\mathrm dx&=e^x+C,\\
\int\sin x\,\mathrm dx&=-\cos x+C, & \int\cos x\,\mathrm dx&=\sin x+C,\\
\int\frac{\mathrm dx}{a^2+x^2}&=\frac1a\arctan\frac xa+C\quad(a>0),\\
\int\frac{\mathrm dx}{\sqrt{a^2-x^2}}&=\arcsin\frac xa+C\quad(a>0,\ |x|<a).
\end{align*}
\subsection{常用不定积分表与递推公式}\index{常用不定积分表与递推公式|hyperpage}
以下 $a>0$，$C$ 为任意常数，公式在被积函数有定义的连通区间内使用。
\begingroup
\allowdisplaybreaks[1]
\begin{align*}
\int a^x\,\mathrm dx&=\frac{a^x}{\ln a}+C\quad(a\ne1),\\
\int\tan x\,\mathrm dx&=-\ln|\cos x|+C,\\
\int\cot x\,\mathrm dx&=\ln|\sin x|+C,\\
\int\sec^2x\,\mathrm dx&=\tan x+C,\\
\int\csc^2x\,\mathrm dx&=-\cot x+C,\\
\int\sec x\tan x\,\mathrm dx&=\sec x+C,\\
\int\csc x\cot x\,\mathrm dx&=-\csc x+C,\\
\int\sec x\,\mathrm dx&=\ln|\sec x+\tan x|+C,\\
\int\csc x\,\mathrm dx&=\ln|\csc x-\cot x|+C,\\
\int\frac{\mathrm dx}{x^2-a^2}&=\frac1{2a}\ln\left|\frac{x-a}{x+a}\right|+C,\\
\int\frac{\mathrm dx}{\sqrt{x^2+a^2}}&=\ln\left|x+\sqrt{x^2+a^2}\right|+C,\\
\int\frac{\mathrm dx}{\sqrt{x^2-a^2}}&=\ln\left|x+\sqrt{x^2-a^2}\right|+C\quad(|x|>a),\\
\int\sqrt{a^2-x^2}\,\mathrm dx&=\frac x2\sqrt{a^2-x^2}+\frac{a^2}2\arcsin(x/a)+C\quad(|x|<a),\\
\int\sinh x\,\mathrm dx&=\cosh x+C,\\
\int\cosh x\,\mathrm dx&=\sinh x+C.
\end{align*}
\endgroup
令 $I_n=\int\sin^n x\,\mathrm dx$，$J_n=\int\cos^n x\,\mathrm dx$，对整数 $n\ge2$，在相差常数的意义下
\[
I_n=-\frac{\sin^{n-1}x\cos x}{n}+\frac{n-1}{n}I_{n-2},\qquad
J_n=\frac{\cos^{n-1}x\sin x}{n}+\frac{n-1}{n}J_{n-2}.
\]
\subsection{定积分的概念与性质}\index{定积分的概念与性质|hyperpage}
对分割 $a=x_0<\cdots<x_n=b$，$\xi_i\in[x_{i-1},x_i]$，
\[
\int_a^b f(x)\,\mathrm dx=\lim_{\max\Delta x_i\to0}\sum_{i=1}^n f(\xi_i)\Delta x_i.
\]
可积函数满足线性、区间可加性与保序性，且
\[
\left|\int_a^b f\right|\le\int_a^b|f|,\quad
\int_a^b f=\int_a^c f+\int_c^b f,\quad
\int_b^a f=-\int_a^b f.
\]
当 $f$ 连续时，$\frac{\mathrm d}{\mathrm dx}\int_a^x f(t)\,\mathrm dt=f(x)$；若 $F'=f$，则
\[
\int_a^b f(x)\,\mathrm dx=F(b)-F(a),\qquad
\int_a^b f(x)\,\mathrm dx=f(\xi)(b-a),\quad \xi\in[a,b].
\]
\subsection{Darboux上下和与Riemann可积性}\index{Darboux上下和与Riemann可积性|hyperpage}
设实函数 $f$ 在 $[a,b]$ 上有界，$P:a=x_0<\cdots<x_n=b$。令
\[
M_i=\sup_{[x_{i-1},x_i]}f,\quad
m_i=\inf_{[x_{i-1},x_i]}f,\quad\Delta x_i=x_i-x_{i-1},
\]
\[
U(f,P)=\sum_{i=1}^nM_i\Delta x_i,\qquad
L(f,P)=\sum_{i=1}^nm_i\Delta x_i.
\]
加细分割使上和不增、下和不减。Darboux 判别为
\[
\begin{aligned}
f\text{ 可积}
&\iff\inf_PU(f,P)=\sup_PL(f,P)\\
&\iff\forall\varepsilon>0,\ \exists P,\quad
\sum_{i=1}^n(M_i-m_i)\Delta x_i<\varepsilon.
\end{aligned}
\]
闭区间上的连续函数、单调函数均可积；仅有有限个间断点的有界函数也可积，特别地，具有有限分段且各段连续延拓至端点的函数可积。
\subsection{加权积分中值定理与积分不等式}\index{加权积分中值定理与积分不等式|hyperpage}
设 $a<b$。若 $f$ 连续，$g$ 可积且不变号，则存在 $\xi\in[a,b]$，使
\[
\int_a^bf(x)g(x)\,\mathrm dx=f(\xi)\int_a^bg(x)\,\mathrm dx.
\]
第二中值定理：若实函数 $f$ 单调、$g$ 连续，则存在 $\xi\in[a,b]$，使
\[
\int_a^bf(x)g(x)\,\mathrm dx
=f(a)\int_a^\xi g(x)\,\mathrm dx
+f(b)\int_\xi^bg(x)\,\mathrm dx.
\]
对实值 Riemann 可积函数 $f,g$，Cauchy--Schwarz 不等式及常用估计为
\[
\left|\int_a^bfg\,\mathrm dx\right|^2
\le\left(\int_a^bf^2\,\mathrm dx\right)
\left(\int_a^bg^2\,\mathrm dx\right),
\]
\[
\left|\int_a^bf\,\mathrm dx\right|
\le\int_a^b|f|\,\mathrm dx
\le(b-a)\sup_{[a,b]}|f|.
\]
\subsection{积分型H\"older与Minkowski不等式}\index{积分型Hölder与Minkowski不等式|hyperpage}
本小节取 $a<b$，实函数 $f,g\in C[a,b]$，并记
\[
\|f\|_p=\left(\int_a^b|f(x)|^p\,\mathrm dx\right)^{1/p}
\quad(1\le p<\infty),\qquad
\|f\|_\infty=\max_{[a,b]}|f|.
\]
若 $p,q>1$ 且 $1/p+1/q=1$，则 H\"older 不等式为
\[
\int_a^b|f(x)g(x)|\,\mathrm dx\le\|f\|_p\|g\|_q.
\]
当两个范数均非零时，等号成立当且仅当存在 $c>0$，使 $|f(x)|^p=c|g(x)|^q$ 对所有 $x\in[a,b]$ 成立；任一函数恒零时也取等号。
对 $1\le p<\infty$，Minkowski 不等式为
\[
\|f+g\|_p\le\|f\|_p+\|g\|_p.
\]
当 $p>1$ 且两函数均非零时，等号成立当且仅当 $f=cg$，其中 $c>0$；$p=1$ 时，等号成立当且仅当 $f(x)g(x)\ge0$ 对所有 $x$ 成立。常用范数比较为
\[
\|f\|_p\le(b-a)^{1/p-1/q}\|f\|_q
\quad(1\le p<q\le\infty),
\]
约定 $1/\infty=0$；非零函数取等号当且仅当 $|f|$ 为常数。测度空间上的推广见第\ref{chap:real-analysis}章。
\par\pagebreak[0]\begin{samepage}
\subsection{积分型Jensen与Hermite--Hadamard不等式}\index{积分型Jensen与Hermite-Hadamard不等式|hyperpage}
设 $\phi$ 为区间 $I$ 上的连续凸函数，$h,w\in C[a,b]$，$h([a,b])\subset I$，$w\ge0$，且 $W=\int_a^bw(x)\,\mathrm dx>0$。则
\[
\phi\!\left(\frac1W\int_a^bh(x)w(x)\,\mathrm dx\right)
\le\frac1W\int_a^b\phi(h(x))w(x)\,\mathrm dx.
\]
若 $\phi$ 严格凸，等号成立当且仅当 $h$ 在 $w>0$ 的集合上取同一常数值。
对 $f\in C[a,b]$ 且凸，Hermite--Hadamard 不等式为
\[
f\!\left(\frac{a+b}2\right)
\le\frac1{b-a}\int_a^bf(x)\,\mathrm dx
\le\frac{f(a)+f(b)}2.
\]
其中任一个不等式取等号当且仅当 $f$ 在 $[a,b]$ 上为仿射函数；严格凸时两侧均为严格不等式。凹函数对应不等号反向，凸性的其他等价条件见第\ref{chap:convex-analysis}章。
\par\end{samepage}\pagebreak[0]
\subsection{积分的计算}\index{积分的计算|hyperpage}
换元与分部积分（函数满足相应可微、可积条件）：
\[
\int_a^b f(\varphi(t))\varphi'(t)\,\mathrm dt
=\int_{\varphi(a)}^{\varphi(b)}f(x)\,\mathrm dx,\qquad
\int_a^b u v'\,\mathrm dx=[uv]_a^b-\int_a^b u'v\,\mathrm dx.
\]
对称区间上，奇函数积分为零，偶函数积分为 $2\int_0^a f(x)\,\mathrm dx$。
\[
\int_a^b f(x)\,\mathrm dx=\int_a^b f(a+b-x)\,\mathrm dx,\qquad
V=\pi\int_a^b[R(x)^2-r(x)^2]\,\mathrm dx.
\]
后一式为绕 $x$ 轴旋转的垫片体积，$0\le r\le R$。
\subsection{弧长、旋转曲面面积与常用定积分}\index{弧长、旋转曲面面积与常用定积分|hyperpage}
对分段 $C^1$ 曲线，
\[
\ell=\int_a^b\sqrt{1+y'(x)^2}\,\mathrm dx,\qquad
\ell=\int_\alpha^\beta\sqrt{x'(t)^2+y'(t)^2}\,\mathrm dt.
\]
极坐标曲线 $r=r(\theta)\ge0$ 的弧长为 $\int_\alpha^\beta\sqrt{r^2+(r')^2}\,\mathrm d\theta$；简单扇形区域面积为 $\frac12\int_\alpha^\beta r^2\,\mathrm d\theta$（无重复覆盖）。
曲线 $y=f(x)\ge0$ 绕 $x$ 轴旋转的曲面面积为
\[
A=2\pi\int_a^b f(x)\sqrt{1+f'(x)^2}\,\mathrm dx.
\]
\begin{align*}
\int_{-\infty}^{\infty}e^{-ax^2}\,\mathrm dx&=\sqrt{\pi/a}\quad(a>0),\\
K_n:=\int_0^{\pi/2}\sin^n x\,\mathrm dx
&=\int_0^{\pi/2}\cos^n x\,\mathrm dx,\\
K_n&=\frac{n-1}{n}K_{n-2}\quad(n\ge2),\qquad K_0=\pi/2,\quad K_1=1.
\end{align*}
\subsection{Gamma函数与Beta积分}\index{Gamma函数与Beta积分|hyperpage}
\begin{gather*}
\Gamma(s)=\int_0^\infty t^{s-1}e^{-t}\,\mathrm dt\quad(s>0),\\
\Gamma(s+1)=s\Gamma(s),\quad\Gamma(n+1)=n!,\quad\Gamma(1/2)=\sqrt\pi,\\
B(p,q)=\int_0^1t^{p-1}(1-t)^{q-1}\,\mathrm dt
=\frac{\Gamma(p)\Gamma(q)}{\Gamma(p+q)}\quad(p,q>0),\\
\int_0^{\pi/2}\sin^{p-1}x\cos^{q-1}x\,\mathrm dx=\tfrac12B(p/2,q/2).
\end{gather*}
\subsection{Gamma、Beta的恒等式与半整数值}\index{Gamma、Beta的恒等式与半整数值|hyperpage}
% Reference: https://dlmf.nist.gov/5.5 and https://dlmf.nist.gov/5.12
在实参数范围内，反射公式与 Legendre 倍增公式为
\[
\Gamma(x)\Gamma(1-x)=\frac\pi{\sin\pi x}\quad(0<x<1),
\]
\[
\Gamma(x)\Gamma\!\left(x+\frac12\right)
=2^{1-2x}\sqrt\pi\,\Gamma(2x)\quad(x>0).
\]
对整数 $n\ge0$，
\[
\Gamma\!\left(n+\frac12\right)
=\frac{(2n)!}{4^nn!}\sqrt\pi.
\]
Beta 函数的对称性、递推及另一种积分表示为（$p,q>0$）
\begin{gather*}
B(p,q)=B(q,p),\qquad
B(p+1,q)=\frac p{p+q}B(p,q),\\
B(p,q+1)=\frac q{p+q}B(p,q),\qquad
B(p,q)=\int_0^\infty\frac{t^{p-1}}{(1+t)^{p+q}}\,\mathrm dt.
\end{gather*}
\subsection{Wallis乘积与Stirling渐近公式}\index{Wallis乘积与Stirling渐近公式|hyperpage}
对 $K_n=\int_0^{\pi/2}\sin^nx\,\mathrm dx$，由递推式得到
\[
K_{2n}=\frac{(2n-1)!!}{(2n)!!}\frac\pi2,\qquad
K_{2n+1}=\frac{(2n)!!}{(2n+1)!!}\quad(n\ge0),
\]
约定 $(-1)!!=0!!=1$。Wallis 乘积为
\[
\frac\pi2=\prod_{n=1}^\infty\frac{(2n)^2}{(2n-1)(2n+1)}.
\]
% Reference: https://dlmf.nist.gov/5.11
当 $n\to\infty$（正整数）时，Stirling 公式及其修正为
\[
n!=\sqrt{2\pi n}\left(\frac ne\right)^n
\left(1+\frac1{12n}+\frac1{288n^2}+O(n^{-3})\right),
\]
\[
\ln(n!)=\left(n+\frac12\right)\ln n-n+\frac12\ln(2\pi)
+\frac1{12n}-\frac1{360n^3}+O(n^{-5}).
\]
相应地，当实数 $x\to+\infty$ 时，
\[
\Gamma(x)\sim\sqrt{2\pi}\,x^{x-1/2}e^{-x}.
\]
这些是渐近公式；修正项按指定阶截断使用，不将其直接视为收敛级数。
\subsection{反常积分的判别}\index{反常积分的判别|hyperpage}
\[
\int_a^\infty f=\lim_{R\to\infty}\int_a^R f,\qquad
\int_a^b f=\lim_{\varepsilon\downarrow0}\int_{a+\varepsilon}^b f
\quad\text{（$a$ 为瑕点）}.
\]
多个瑕点或无穷端点必须分别收敛。
\[
\int_1^\infty\frac{\mathrm dx}{x^p}\text{ 收敛}\iff p>1,\qquad
\int_0^1\frac{\mathrm dx}{x^p}\text{ 收敛}\iff p<1.
\]
比较判别：$0\le f\le g$，$\int g$ 收敛则 $\int f$ 收敛；$\int f$ 发散则 $\int g$ 发散。
若非负函数 $f/g\to c\in(0,\infty)$，则同敛散；$\int|f|<\infty$ 蕴含 $\int f$ 收敛。
\subsection{反常积分的Cauchy、Dirichlet与Abel判别}\index{反常积分的Cauchy、Dirichlet与Abel判别|hyperpage}
设 $f$ 在 $[a,+\infty)$ 的每个紧区间上可积。Cauchy 准则为
\[
\int_a^\infty f\text{ 收敛}\iff
\forall\varepsilon>0,\ \exists R>a,\ \forall B>A\ge R,\quad
\left|\int_A^B f(x)\,\mathrm dx\right|<\varepsilon.
\]
Dirichlet 判别：若原函数型积分 $F(R)=\int_a^Rf(x)\,\mathrm dx$ 有界，$g\in C^1([a,\infty))$ 单调且趋于零，则 $\int_a^\infty fg$ 收敛。
Abel 判别：若 $\int_a^\infty f$ 收敛，$g\in C^1([a,\infty))$ 单调且有界，则 $\int_a^\infty fg$ 收敛。二者均可由
\[
\int_A^B f(x)g(x)\,\mathrm dx
=[F(x)g(x)]_A^B-\int_A^BF(x)g'(x)\,\mathrm dx
\]
得到；有限瑕点的对应结论可经换元转化后使用。
\section{有界变差与Riemann--Stieltjes积分}
\subsection{总变差与Jordan分解}\index{总变差与Jordan分解|hyperpage}
% Reference: https://staff.fnwi.uva.nl/p.j.c.spreij/onderwijs/master/bv.pdf
对实函数 $f:[a,b]\to\mathbb R$，总变差定义为
\[
V_a^b(f)=\sup_{P}\sum_{i=1}^n|f(x_i)-f(x_{i-1})|,
\quad P:a=x_0<\cdots<x_n=b.
\]
若 $V_a^b(f)<\infty$，记 $f\in BV[a,b]$。有界变差函数有界，且在内点具有有限左右极限；本节记 $\|f\|_\infty=\sup_{[a,b]}|f|$。基本性质为
\begin{gather*}
V_a^b(f)=V_a^c(f)+V_c^b(f)\quad(a\le c\le b),\\
V_a^b(\lambda f)=|\lambda|V_a^b(f),\qquad
V_a^b(f+g)\le V_a^b(f)+V_a^b(g),\\
V_a^b(fg)\le\|f\|_\infty V_a^b(g)+\|g\|_\infty V_a^b(f).
\end{gather*}
单调函数满足 $V_a^b(f)=|f(b)-f(a)|$；若 $f\in C^1[a,b]$，则
\[
V_a^b(f)=\int_a^b|f'(x)|\,\mathrm dx.
\]
\par\pagebreak[0]\begin{samepage}
Jordan 分解：令 $v(x)=V_a^x(f)$，则
\[
p(x)=\frac{v(x)+f(x)-f(a)}2,\qquad
q(x)=\frac{v(x)-f(x)+f(a)}2
\]
均单调不减，且 $p(a)=q(a)=0$、$f(x)=f(a)+p(x)-q(x)$。因此 $f\in BV$ 当且仅当它可表示为两个单调不减函数之差。
\par\end{samepage}\pagebreak[0]
\subsection{Stieltjes积分的定义、存在与估计}\index{Stieltjes积分的定义、存在与估计|hyperpage}
对带标记的分割 $P$，$\xi_i\in[x_{i-1},x_i]$，定义
\[
\int_a^bf\,\mathrm d\alpha
=\lim_{\|P\|\to0}\sum_{i=1}^nf(\xi_i)
\bigl[\alpha(x_i)-\alpha(x_{i-1})\bigr],
\qquad\|P\|=\max_i(x_i-x_{i-1}),
\]
这里要求极限与标记点和分割的选取无关。若 $f\in C[a,b]$、$\alpha\in BV[a,b]$，则积分存在，并且
\[
\left|\int_a^bf\,\mathrm d\alpha\right|
\le\|f\|_\infty V_a^b(\alpha).
\]
若 $\alpha\in C^1[a,b]$，则化为普通积分：
\[
\int_a^bf(x)\,\mathrm d\alpha(x)
=\int_a^bf(x)\alpha'(x)\,\mathrm dx.
\]
若 $\alpha$ 单调不减，则积分保序；对连续实函数 $f$，还存在 $\xi\in[a,b]$ 使
\[
\int_a^bf\,\mathrm d\alpha=f(\xi)\,[\alpha(b)-\alpha(a)].
\]
\subsection{Stieltjes分部积分与一致极限}\index{Stieltjes分部积分与一致极限|hyperpage}
若 $f,g\in C[a,b]\cap BV[a,b]$，则
\[
\int_a^bf\,\mathrm dg+\int_a^bg\,\mathrm df
=f(b)g(b)-f(a)g(a).
\]
连续性条件使上述形式无需跳跃项。若 $f_n,f\in C[a,b]$，$f_n\rightrightarrows f$，且固定的 $\alpha\in BV[a,b]$，则
\[
\left|\int_a^b(f_n-f)\,\mathrm d\alpha\right|
\le\|f_n-f\|_\infty V_a^b(\alpha)\longrightarrow0.
\]
若连续函数项级数 $\sum u_n$ 在 $[a,b]$ 上一致收敛，则可逐项作 Stieltjes 积分：
\[
\int_a^b\sum_{n=1}^\infty u_n\,\mathrm d\alpha
=\sum_{n=1}^\infty\int_a^bu_n\,\mathrm d\alpha.
\]
\section{含参积分}
\subsection{有限区间上的连续性与Leibniz公式}\index{有限区间上的连续性与Leibniz公式|hyperpage}
设 $I$ 为参数的开区间，$f(x,t)$ 在 $[a,b]\times I$ 上连续，则
\[
F(t)=\int_a^bf(x,t)\,\mathrm dx
\]
在 $I$ 上连续。若 $\partial_t f$ 也连续，则
\[
F'(t)=\int_a^b\partial_t f(x,t)\,\mathrm dx.
\]
若 $\alpha,\beta\in C^1(I)$，$f$ 与 $\partial_t f$ 在变限所经过区域的邻域内连续，则
\[
\begin{aligned}
\frac{\mathrm d}{\mathrm dt}\int_{\alpha(t)}^{\beta(t)}f(x,t)\,\mathrm dx
={}&f(\beta(t),t)\beta'(t)-f(\alpha(t),t)\alpha'(t)\\
&+\int_{\alpha(t)}^{\beta(t)}\partial_t f(x,t)\,\mathrm dx.
\end{aligned}
\]
\subsection{反常含参积分的一致收敛}\index{反常含参积分的一致收敛|hyperpage}
\par\pagebreak[0]\begin{samepage}
设 $f(\cdot,t)$ 在 $[a,\infty)$ 上局部可积，$t\in E$。定义
\[
\int_a^\infty f(x,t)\,\mathrm dx\text{ 于 }E\text{ 一致收敛}
\iff
\sup_{t\in E}\left|\int_R^\infty f(x,t)\,\mathrm dx\right|\xrightarrow[R\to\infty]{}0,
\]
此式右侧要求各参数下的积分均已存在。
\par\end{samepage}\pagebreak[0]
等价的统一 Cauchy 准则为
\[
\forall\varepsilon>0,\ \exists R>a,\ \forall B>A\ge R,\ \forall t\in E,
\quad\left|\int_A^B f(x,t)\,\mathrm dx\right|<\varepsilon.
\]
Weierstrass 判别：若 $|f(x,t)|\le g(x)$ 对所有 $t\in E$ 成立，且 $g\ge0$、$\int_a^\infty g<\infty$，则该含参积分一致收敛。
一致 Dirichlet 判别：若 $\left|\int_a^Rf(x,t)\,\mathrm dx\right|\le M$ 对 $R\ge a,t\in E$ 一致成立，$g(\cdot,t)$ 为 $C^1$ 单调函数，且 $\sup_{t\in E}|g(x,t)|\to0$，则 $\int_a^\infty f(x,t)g(x,t)\,\mathrm dx$ 一致收敛。
\subsection{反常含参积分的极限、求导与积分交换}\index{反常含参积分的极限、求导与积分交换|hyperpage}
设 $I$ 为开区间，$f$ 在 $[a,\infty)\times I$ 上连续，且 $\int_a^\infty f(x,t)\,\mathrm dx$ 在每个紧参数区间上一致收敛，则 $F(t)$ 连续，并有
\[
\lim_{t\to t_0}\int_a^\infty f(x,t)\,\mathrm dx
=\int_a^\infty f(x,t_0)\,\mathrm dx\quad(t_0\in I).
\]
若 $\partial_t f$ 连续，且 $\int_a^\infty\partial_t f(x,t)\,\mathrm dx$ 也在每个紧参数区间上一致收敛，则
\[
F'(t)=\int_a^\infty\partial_t f(x,t)\,\mathrm dx.
\]
一组便于核验的充分条件是：对每个紧区间 $J\subset I$，存在可积的非负函数 $g_J,h_J$，使
\[
|f(x,t)|\le g_J(x),\quad |\partial_t f(x,t)|\le h_J(x)
\quad(x\ge a,\ t\in J).
\]
在上述 $f$ 连续、积分局部一致收敛的条件下，对 $[c,d]\subset I$ 还可交换积分次序：
\[
\int_c^d\!\left(\int_a^\infty f(x,t)\,\mathrm dx\right)\mathrm dt
=\int_a^\infty\!\left(\int_c^df(x,t)\,\mathrm dt\right)\mathrm dx.
\]
\subsection{Frullani积分公式}\index{Frullani积分公式|hyperpage}
设 $f\in C(0,\infty)$，并存在有限极限
\[
f(0+)=\lim_{x\downarrow0}f(x),\qquad
f(\infty)=\lim_{x\to\infty}f(x).
\]
对 $a,b>0$，反常积分收敛且
\[
\int_0^\infty\frac{f(ax)-f(bx)}x\,\mathrm dx
=[f(0+)-f(\infty)]\ln\frac ba.
\]
条件保证两端分别收敛；计算中应保留差式作为整体，不能拆成可能各自发散的两个积分。
\subsection{指数振荡积分与Dirichlet积分}\index{指数振荡积分与Dirichlet积分|hyperpage}
对 $s>0$、$b\in\mathbb R$，有绝对收敛公式
\begin{align*}
\int_0^\infty e^{-sx}\cos(bx)\,\mathrm dx&=\frac{s}{s^2+b^2},\\
\int_0^\infty e^{-sx}\sin(bx)\,\mathrm dx&=\frac{b}{s^2+b^2},\\
\int_0^\infty e^{-sx}\frac{\sin(bx)}x\,\mathrm dx
&=\arctan\frac bs.
\end{align*}
Dirichlet 积分为
\[
\int_0^\infty\frac{\sin(bx)}x\,\mathrm dx
=\frac\pi2\operatorname{sgn}(b)\quad(b\ne0).
\]
$b\ne0$ 时该积分条件收敛；$b=0$ 时被积函数取零，积分为零。上式也等于带衰减积分在 $s\downarrow0$ 时的极限，此处的极限交换应由反常积分的统一尾部估计保证。
\section{向量代数与空间解析几何}
\subsection{向量的概念与线性运算}\index{向量的概念与线性运算|hyperpage}
设 $\bm a=(a_1,a_2,a_3)$，$\bm b=(b_1,b_2,b_3)$。
\[
\bm a\pm\bm b=(a_1\pm b_1,a_2\pm b_2,a_3\pm b_3),\quad
\lambda\bm a=(\lambda a_1,\lambda a_2,\lambda a_3),\quad
|\bm a|=\sqrt{a_1^2+a_2^2+a_3^2}.
\]
\[
\overrightarrow{AB}=\bm r_B-\bm r_A,\qquad
\bm e=\bm a/|\bm a|\ (\bm a\ne0),\qquad
\bm r=\bm r_0+t\bm v\ (\bm v\ne0)\quad\text{（直线）}.
\]
\subsection{向量的运算与线性运算}\index{向量的运算与线性运算|hyperpage}
\[
\bm a\cdot\bm b=\sum_{i=1}^3a_ib_i=|\bm a||\bm b|\cos\theta,\qquad
\bm a\times\bm b=\begin{vmatrix}\bm i&\bm j&\bm k\\a_1&a_2&a_3\\b_1&b_2&b_3\end{vmatrix}.
\]
\[
|\bm a\times\bm b|=|\bm a||\bm b|\sin\theta,\quad
V_{\text{平行六面体}}=|\bm a\cdot(\bm b\times\bm c)|.
\]
平面方程及点到平面距离为
\[
Ax+By+Cz+D=0,\qquad
 d=\frac{|Ax_0+By_0+Cz_0+D|}{\sqrt{A^2+B^2+C^2}}\quad(A^2+B^2+C^2>0).
\]
点到直线 $\bm r=\bm r_0+t\bm v$ 的距离为 $| (\bm r_P-\bm r_0)\times\bm v|/|\bm v|$。
\section{多元函数微分学}
\subsection{多元函数的概念，多元函数的集合意义，二元函数的极限和连续的概念}\index{多元函数的概念，多元函数的集合意义，二元函数的极限和连续的概念|hyperpage}
$z=f(x,y)$ 的定义域 $D\subseteq\mathbb R^2$，图像为 $\{(x,y,f(x,y)):(x,y)\in D\}$。
\[
\begin{gathered}
\lim_{(x,y)\to(a,b)}f(x,y)=L\iff\forall\varepsilon>0,\ \exists\delta>0:\\
0<\sqrt{(x-a)^2+(y-b)^2}<\delta\Longrightarrow|f(x,y)-L|<\varepsilon.
\end{gathered}
\]
极限沿定义域内任意路径均须相同；连续要求该极限等于 $f(a,b)$。可微的定义为
\[
\Delta f=f_x\Delta x+f_y\Delta y+o\!\left(\sqrt{\Delta x^2+\Delta y^2}\right),\qquad
\mathrm df=f_x\,\mathrm dx+f_y\,\mathrm dy.
\]
偏导数在邻域内连续是可微的充分条件，可微蕴含连续。
\subsection{全导数、Jacobian矩阵与复合映射}\index{全导数、Jacobian矩阵与复合映射|hyperpage}
\par\pagebreak[0]\begin{samepage}
设 $U\subset\mathbb R^n$ 为开集，$f:U\to\mathbb R^m$。若存在线性映射 $Df(\bm a)$，使
\[
f(\bm a+\bm h)=f(\bm a)+Df(\bm a)\bm h+o(\|\bm h\|),
\]
则称 $f$ 在 $\bm a$ 可微。
\par\end{samepage}\pagebreak[0]
其标准基下的矩阵为
\[
J_f(\bm a)=\left(\frac{\partial f_i}{\partial x_j}(\bm a)\right)
_{\substack{1\le i\le m\\1\le j\le n}},\qquad
\mathrm d\bm y=J_f(\bm x)\,\mathrm d\bm x.
\]
偏导数连续是可微的充分条件；仅有各偏导数存在不能保证可微。
\par\pagebreak[0]\begin{samepage}
若 $g:\mathbb R^m\supset V\to\mathbb R^p$ 在 $f(\bm a)$ 可微，则
\[
D(g\circ f)(\bm a)=Dg(f(\bm a))Df(\bm a),\qquad
J_{g\circ f}(\bm a)=J_g(f(\bm a))J_f(\bm a).
\]
\par\end{samepage}\pagebreak[0]
\subsection{混合偏导交换与多元有限增量公式}\index{混合偏导交换与多元有限增量公式|hyperpage}
Schwarz 定理：若 $f\in C^2(U)$，$U\subset\mathbb R^n$ 为开集，则
\[
\frac{\partial^2f}{\partial x_i\partial x_j}
=\frac{\partial^2f}{\partial x_j\partial x_i}\quad(1\le i,j\le n).
\]
若 $f\in C^k$，则总阶数不超过 $k$ 的混合偏导可任意交换求导次序。仅有混合偏导存在不足以保证交换。
设 $f:U\to\mathbb R^m$ 为 $C^1$，线段 $[\bm a,\bm a+\bm h]\subset U$，则
\[
f(\bm a+\bm h)-f(\bm a)
=\int_0^1Df(\bm a+t\bm h)\bm h\,\mathrm dt.
\]
取欧氏范数及其诱导的算子范数，有
\[
\|f(\bm a+\bm h)-f(\bm a)\|
\le\sup_{0\le t\le1}\|Df(\bm a+t\bm h)\|_{\mathrm{op}}\,\|\bm h\|.
\]
因此在凸开集 $U$ 上，$\sup_U\|Df\|_{\mathrm{op}}\le M$ 蕴含 $f$ 为 Lipschitz 映射。标量值情形还存在 $t_*\in(0,1)$，使
\[
f(\bm a+\bm h)-f(\bm a)=\nabla f(\bm a+t_*\bm h)\cdot\bm h.
\]
向量值情形一般不能要求所有分量共用同一个中间点，应使用积分形式。
\subsection{反函数定理与一般隐函数定理}\index{反函数定理与一般隐函数定理|hyperpage}
若 $f:U\subset\mathbb R^n\to\mathbb R^n$ 为 $C^1$ 映射，且 $\det J_f(\bm a)\ne0$，则 $f$ 在 $\bm a$ 附近有 $C^1$ 局部反函数，并且
\[
D(f^{-1})(f(\bm x))=[Df(\bm x)]^{-1},\qquad
\det J_{f^{-1}}(f(\bm x))=\frac1{\det J_f(\bm x)}.
\]
此结论只保证局部可逆，不能单凭 Jacobian 处处非零断定全局一一对应。
设 $F:U\subset\mathbb R^n\times\mathbb R^m\to\mathbb R^m$ 为 $C^1$，且
\[
F(\bm a,\bm b)=\bm0,\qquad \det D_{\bm y}F(\bm a,\bm b)\ne0.
\]
则在该点附近，方程 $F(\bm x,\bm y)=\bm0$ 唯一确定 $C^1$ 函数 $\bm y=\varphi(\bm x)$，满足
\[
\varphi(\bm a)=\bm b,\qquad
D\varphi(\bm x)=-[D_{\bm y}F(\bm x,\varphi(\bm x))]^{-1}
D_{\bm x}F(\bm x,\varphi(\bm x)).
\]
若原映射为 $C^k$（$k\ge1$），上述局部反函数或隐函数也为 $C^k$。
\subsection{隐函数求导法}\index{隐函数求导法|hyperpage}
若 $F(x,y)=0$ 且 $F\in C^1$、$F_y\ne0$，则局部有
\[
y'=-\frac{F_x}{F_y}.
\]
若 $F\in C^2$，则 $y''=-[F_{xx}+2F_{xy}y'+F_{yy}(y')^2]/F_y$。
对 $F(x,y,z)=0$，在 $F_z\ne0$ 且 $F\in C^1$ 时
\[
z_x=-F_x/F_z,\qquad z_y=-F_y/F_z.
\]
复合函数 $z=f(u(x,y),v(x,y))$ 满足 $z_x=f_u u_x+f_v v_x$，$z_y=f_u u_y+f_v v_y$。
\subsection{几何公式}\index{几何公式|hyperpage}
曲面 $z=f(x,y)$ 在 $(x_0,y_0,z_0)$ 的切平面为
\[
z-z_0=f_x(x_0,y_0)(x-x_0)+f_y(x_0,y_0)(y-y_0).
\]
隐曲面 $F(x,y,z)=0$ 在 $P$ 处若 $\nabla F(P)\ne0$，则
\[
\text{切平面： }\nabla F(P)\cdot(\bm r-\bm r_P)=0,\qquad
\text{法线： }\bm r=\bm r_P+t\nabla F(P).
\]
空间正则曲线 $\bm r(t)$ 的切线为 $\bm r=\bm r(t_0)+s\bm r'(t_0)$。
\subsection{多元Taylor公式与Hessian矩阵}\index{多元Taylor公式与Hessian矩阵|hyperpage}
设 $f$ 在 $\bm a$ 附近二阶连续可微，Hessian 矩阵为
\[
H_f=\left(\frac{\partial^2f}{\partial x_i\partial x_j}\right)_{i,j=1}^n.
\]
当 $\bm h\to\bm0$ 时，
\[
f(\bm a+\bm h)=f(\bm a)+\nabla f(\bm a)^{\mathsf T}\bm h
+\tfrac12\bm h^{\mathsf T}H_f(\bm a)\bm h+o(\|\bm h\|^2).
\]
驻点处 $H_f$ 正定则为严格极小点，负定则为严格极大点，不定则为鞍点；半定时此判别不足。
约束极值中取 $L=f-\sum_i\lambda_i g_i$，约束梯度线性无关时，在满足 $\nabla g_i^{\mathsf T}\bm h=0$ 的非零切向量上，若 $\bm h^{\mathsf T}H_{\bm x}L\bm h>0$，则为严格条件极小（负号对应极大）。此处假定目标与约束均为 $C^2$。
\subsection{多重指标Taylor公式与齐次函数}\index{多重指标Taylor公式与齐次函数|hyperpage}
对 $\alpha=(\alpha_1,\ldots,\alpha_n)\in\mathbb N_0^n$，记
\[
|\alpha|=\sum_{i=1}^n\alpha_i,\quad
\alpha!=\prod_{i=1}^n\alpha_i!,\quad
\bm h^\alpha=\prod_{i=1}^nh_i^{\alpha_i},\quad
D^\alpha f=\frac{\partial^{|\alpha|}f}{\partial x_1^{\alpha_1}\cdots\partial x_n^{\alpha_n}}.
\]
若 $f\in C^m$ 于 $\bm a$ 的邻域，$m\ge1$，则
\[
f(\bm a+\bm h)=\sum_{|\alpha|\le m}
\frac{D^\alpha f(\bm a)}{\alpha!}\bm h^\alpha+o(\|\bm h\|^m).
\]
若 $f\in C^{m+1}$ 且连接 $\bm a,\bm a+\bm h$ 的线段位于定义域内，积分型余项为
\[
R_m(\bm h)=(m+1)\sum_{|\alpha|=m+1}\frac{\bm h^\alpha}{\alpha!}
\int_0^1(1-t)^mD^\alpha f(\bm a+t\bm h)\,\mathrm dt.
\]
若线段上所有 $m+1$ 阶偏导的绝对值均不超过 $M$，则
\[
|R_m(\bm h)|\le\frac{M}{(m+1)!}
\left(\sum_{i=1}^n|h_i|\right)^{m+1}.
\]
设 $f$ 在开锥域上为 $C^1$，且 $f(t\bm x)=t^kf(\bm x)$（$t>0$），则 Euler 公式为
\[
\bm x\cdot\nabla f(\bm x)=kf(\bm x).
\]
若 $f\in C^2$，进一步有
\[
H_f(\bm x)\bm x=(k-1)\nabla f(\bm x),\qquad
\bm x^{\mathsf T}H_f(\bm x)\bm x=k(k-1)f(\bm x).
\]
\subsection{极值}\index{极值|hyperpage}
内点可微极值的必要条件为 $\nabla f=0$。对二元 $C^2$ 函数在驻点处令
\[
A=f_{xx},\quad B=f_{xy},\quad C=f_{yy},\quad D=AC-B^2.
\]
$D>0,A>0$ 为严格极小；$D>0,A<0$ 为严格极大；$D<0$ 为鞍点；$D=0$ 时需另行判断。
等式约束 $g_i(\bm x)=0$ 的梯度线性无关时，条件极值满足
\[
\nabla f=\sum_{i=1}^m\lambda_i\nabla g_i,\qquad g_i=0.
\]
闭有界区域上连续函数的最值需比较内部驻点与边界候选点。
\subsection{方向导数与梯度}\index{方向导数与梯度|hyperpage}
\par\pagebreak[0]\begin{samepage}
对单位向量 $\bm e$，可微函数有
\[
D_{\bm e}f(P)=\lim_{t\to0}\frac{f(P+t\bm e)-f(P)}t
=\nabla f(P)\cdot\bm e,\qquad \nabla f=(f_{x_1},\ldots,f_{x_n}).
\]
方向导数最大值为 $|\nabla f|$，梯度非零时在梯度方向取得；梯度垂直于正则等值面。
\par\end{samepage}\pagebreak[0]
\section{多元函数积分学}
以下换元要求变换在区域内部一一对应、连续可微且 Jacobian 非零；边界上允许通常的零测例外。积分交换次序要求可积（如连续或绝对可积）。
\subsection{二重积分}\index{二重积分|hyperpage}
\[
\iint_D f\,\mathrm dA=\int_a^b\!\mathrm dx\int_{\varphi_1(x)}^{\varphi_2(x)}f(x,y)\,\mathrm dy,
\quad D=\{(x,y):a\le x\le b,\ \varphi_1(x)\le y\le\varphi_2(x)\}.
\]
\[
\iint_D f(x,y)\,\mathrm dx\,\mathrm dy
=\iint_G f(x(u,v),y(u,v))\left|\frac{\partial(x,y)}{\partial(u,v)}\right|\,\mathrm du\,\mathrm dv.
\]
极坐标 $x=r\cos\theta,y=r\sin\theta$，$r\ge0$，$\mathrm dA=r\,\mathrm dr\,\mathrm d\theta$。
区域面积 $|D|=\iint_D1\,\mathrm dA$，质量 $M=\iint_D\rho\,\mathrm dA$。
\subsection{三重积分}\index{三重积分|hyperpage}
\[
\iiint_\Omega f\,\mathrm dV=\iint_D\!\mathrm dx\,\mathrm dy
\int_{z_1(x,y)}^{z_2(x,y)}f(x,y,z)\,\mathrm dz.
\]
柱坐标：$x=r\cos\theta,y=r\sin\theta,z=z$，$\mathrm dV=r\,\mathrm dr\,\mathrm d\theta\,\mathrm dz$。
球坐标：$x=r\sin\varphi\cos\theta,y=r\sin\varphi\sin\theta,z=r\cos\varphi$，
\[
\mathrm dV=r^2\sin\varphi\,\mathrm dr\,\mathrm d\varphi\,\mathrm d\theta,
\quad r\ge0,\quad0\le\varphi\le\pi,\quad0\le\theta<2\pi.
\]
一般换元体积元为 $\mathrm dV=|\partial(x,y,z)/\partial(u,v,w)|\,\mathrm du\,\mathrm dv\,\mathrm dw$。
\subsection{高维球体积与球面面积}\index{高维球体积与球面面积|hyperpage}
对整数 $n\ge2$、$R>0$，记 $B_n(R)=\{\bm x\in\mathbb R^n:\|\bm x\|\le R\}$，则
\[
V_n(R)=|B_n(R)|=\frac{\pi^{n/2}}{\Gamma(n/2+1)}R^n,
\]
\[
|\partial B_n(R)|=\frac{\mathrm d}{\mathrm dR}V_n(R)
=\frac{2\pi^{n/2}}{\Gamma(n/2)}R^{n-1}.
\]
单位球面面积记为 $\sigma_{n-1}=2\pi^{n/2}/\Gamma(n/2)$，故 $V_n(R)=\sigma_{n-1}R^n/n$。这里的球面面积为 $n-1$ 维面积。
\subsection{高维极坐标与径向积分}\index{高维极坐标与径向积分|hyperpage}
在 $\mathbb R^n$ 中写 $\bm x=r\bm\omega$，$r>0$，$\bm\omega\in S^{n-1}$。极坐标积分公式为
\[
\int_{\mathbb R^n}f(\bm x)\,\mathrm d\bm x
=\int_0^\infty r^{n-1}
\left(\int_{S^{n-1}}f(r\bm\omega)\,\mathrm dS_{\bm\omega}\right)\mathrm dr.
\]
此处对非负可测函数允许两端为 $+\infty$；有符号函数采用绝对可积条件。在经典多重积分中可先用于有界区域上的连续函数，再按绝对收敛取极限。
特别地，对径向函数 $f(\bm x)=F(\|\bm x\|)$，
\[
\int_{\mathbb R^n}F(\|\bm x\|)\,\mathrm d\bm x
=\sigma_{n-1}\int_0^\infty F(r)r^{n-1}\,\mathrm dr,
\]
\[
\int_{B_n(R)}F(\|\bm x\|)\,\mathrm d\bm x
=\sigma_{n-1}\int_0^RF(r)r^{n-1}\,\mathrm dr.
\]
\subsection{Dirichlet单纯形积分}\index{Dirichlet单纯形积分|hyperpage}
% Reference: https://dlmf.nist.gov/5.14
设 $n\ge1$，标准单纯形为
\[
\Delta_n=\left\{\bm x\in\mathbb R^n:x_i\ge0,\ \sum_{i=1}^nx_i\le1\right\}.
\]
当 $\alpha_0,\alpha_1,\ldots,\alpha_n>0$ 时，积分在边界可能有瑕点，但收敛，并满足
\[
\int_{\Delta_n}\left(1-\sum_{i=1}^nx_i\right)^{\alpha_0-1}
\prod_{i=1}^nx_i^{\alpha_i-1}\,\mathrm d\bm x
=\frac{\prod_{i=0}^n\Gamma(\alpha_i)}
{\Gamma\!\left(\sum_{i=0}^n\alpha_i\right)}.
\]
二维形式为（$p,q,r>0$）
\[
\iint_{\Delta_2}x^{p-1}y^{q-1}(1-x-y)^{r-1}\,\mathrm dx\,\mathrm dy
=\frac{\Gamma(p)\Gamma(q)\Gamma(r)}{\Gamma(p+q+r)}.
\]
相应的单纯形体积为 $|\Delta_n|=1/n!$；当 $a_i>0$ 时，区域 $x_i\ge0$、$\sum_i x_i/a_i\le1$ 的体积为 $\prod_i a_i/n!$。
\subsection{质量、静矩与质心}\index{质量、静矩与质心|hyperpage}
对空间区域 $\Omega$ 上的非负密度 $\rho$，若质量及绝对一阶矩均有限，且 $M>0$，则
\[
M=\iiint_\Omega\rho\,\mathrm dV,\qquad
\bar{\bm r}=\frac1M\iiint_\Omega\bm r\rho\,\mathrm dV.
\]
关于三个坐标平面的静矩及质心坐标为
\begin{gather*}
S_{yz}=\iiint_\Omega x\rho\,\mathrm dV,\quad
S_{xz}=\iiint_\Omega y\rho\,\mathrm dV,\quad
S_{xy}=\iiint_\Omega z\rho\,\mathrm dV,\\
\bar x=S_{yz}/M,\qquad \bar y=S_{xz}/M,\qquad \bar z=S_{xy}/M.
\end{gather*}
对平面薄片 $D$、曲线 $L$ 和曲面 $S$，统一记相应密度为 $\rho$，有
\begin{align*}
M_D&=\iint_D\rho\,\mathrm dA,
&\bar{\bm r}_D&=\frac1{M_D}\iint_D\bm r\rho\,\mathrm dA,\\
M_L&=\int_L\rho\,\mathrm ds,
&\bar{\bm r}_L&=\frac1{M_L}\int_L\bm r\rho\,\mathrm ds,\\
M_S&=\iint_S\rho\,\mathrm dS,
&\bar{\bm r}_S&=\frac1{M_S}\iint_S\bm r\rho\,\mathrm dS.
\end{align*}
均匀密度时，质心等于几何重心；若区域及密度关于某平面反射对称，则质心位于该平面。
\subsection{转动惯量与平行轴定理}\index{转动惯量与平行轴定理|hyperpage}
设质量有限、非零，且二阶矩有限。到转轴 $\ell$ 的距离为 $d(\bm r,\ell)$，则
\[
I_\ell=\int d(\bm r,\ell)^2\,\mathrm dm,
\qquad\mathrm dm=\rho\,\mathrm dV
\]
（薄片、曲线、曲面分别使用相应质量元）。空间物体绕坐标轴的转动惯量为
\begin{align*}
I_x&=\iiint_\Omega(y^2+z^2)\rho\,\mathrm dV,\\
I_y&=\iiint_\Omega(x^2+z^2)\rho\,\mathrm dV,\\
I_z&=\iiint_\Omega(x^2+y^2)\rho\,\mathrm dV.
\end{align*}
若轴 $\ell_C$ 经过质心且与 $\ell$ 平行，两轴间距为 $d$，则
\[
I_\ell=I_{\ell_C}+Md^2.
\]
对位于 $z=0$ 平面的薄片，垂直轴定理为 $I_z=I_x+I_y$，其中三轴经过同一点且 $x,y$ 轴位于薄片平面内。
\subsection{第一类曲线积分}\index{第一类曲线积分|hyperpage}
\par\pagebreak[0]\begin{samepage}
对分段光滑曲线 $\bm r(t)$，$a\le t\le b$，
\[
\int_L f\,\mathrm ds=\int_a^b f(\bm r(t))|\bm r'(t)|\,\mathrm dt,
\quad \mathrm ds=\sqrt{x'^2+y'^2+z'^2}\,\mathrm dt.
\]
曲线长度 $\ell=\int_L1\,\mathrm ds$；第一类曲线积分与方向无关。
\par\end{samepage}\pagebreak[0]
\subsection{第二类曲线积分}\index{第二类曲线积分|hyperpage}
\[
\int_L P\,\mathrm dx+Q\,\mathrm dy+R\,\mathrm dz
=\int_a^b[P(\bm r)x'+Q(\bm r)y'+R(\bm r)z']\,\mathrm dt.
\]
反向时积分变号。若 $\bm F=\nabla\phi$，则 $\int_A^B\bm F\cdot\mathrm d\bm r=\phi(B)-\phi(A)$。
平面区域 $D$ 的分段光滑边界取正向（区域在行进方向左侧），且 $P,Q\in C^1$ 时，Green 公式为
\[
\oint_{\partial D}P\,\mathrm dx+Q\,\mathrm dy
=\iint_D(Q_x-P_y)\,\mathrm dx\,\mathrm dy.
\]
\subsection{第一类曲面积分}\index{第一类曲面积分|hyperpage}
对正则参数曲面 $\bm r(u,v)$，
\[
\iint_S f\,\mathrm dS=\iint_D f(\bm r(u,v))|\bm r_u\times\bm r_v|\,\mathrm du\,\mathrm dv.
\]
若 $z=z(x,y)$，则 $\mathrm dS=\sqrt{1+z_x^2+z_y^2}\,\mathrm dx\,\mathrm dy$。
曲面面积 $|S|=\iint_S1\,\mathrm dS$，积分与定向无关。
\subsection{第二类曲面积分}\index{第二类曲面积分|hyperpage}
对选定单位法向量 $\bm n$，通量为
\[
\iint_S\bm F\cdot\bm n\,\mathrm dS
=\iint_S P\,\mathrm dy\,\mathrm dz+Q\,\mathrm dz\,\mathrm dx+R\,\mathrm dx\,\mathrm dy.
\]
若参数化的 $\bm r_u\times\bm r_v$ 与定向相容，则通量为 $\iint_D\bm F(\bm r)\cdot(\bm r_u\times\bm r_v)\,\mathrm du\,\mathrm dv$。图形曲面 $z=z(x,y)$ 取向上侧时为
\[
\iint_D(-Pz_x-Qz_y+R)\,\mathrm dx\,\mathrm dy.
\]
反向时通量变号。
\subsection{散度与旋度}\index{散度与旋度|hyperpage}
\[
\operatorname{div}\bm F=\nabla\cdot\bm F=P_x+Q_y+R_z,\qquad
\operatorname{curl}\bm F=\nabla\times\bm F=(R_y-Q_z,P_z-R_x,Q_x-P_y).
\]
对 $C^1$ 向量场及分段光滑的适当区域，Gauss 与 Stokes 公式为
\[
\oiint_{\partial\Omega}\bm F\cdot\bm n\,\mathrm dS
=\iiint_\Omega\nabla\cdot\bm F\,\mathrm dV,\qquad
\oint_{\partial S}\bm F\cdot\mathrm d\bm r
=\iint_S(\nabla\times\bm F)\cdot\bm n\,\mathrm dS.
\]
Gauss 公式取外法向；Stokes 公式边界方向与法向遵循右手定则。二阶连续可微时，
\[
\nabla\times(\nabla f)=\bm0,\quad
\nabla\cdot(\nabla\times\bm F)=0,\quad
\Delta f=\nabla\cdot\nabla f=f_{xx}+f_{yy}+f_{zz}.
\]
\subsection{向量微分的乘积公式与二阶恒等式}\index{向量微分的乘积公式与二阶恒等式|hyperpage}
以下在三维直角坐标中使用，$f,g,\bm F,\bm G$ 具有公式所需的一阶或二阶连续导数。
\begin{align*}
\nabla(fg)&=f\nabla g+g\nabla f,\\
\operatorname{div}(f\bm F)&=\nabla f\cdot\bm F+f\,\operatorname{div}\bm F,\\
\operatorname{curl}(f\bm F)&=\nabla f\times\bm F+f\,\operatorname{curl}\bm F,\\
\operatorname{div}(\bm F\times\bm G)
&=\bm G\cdot\operatorname{curl}\bm F-\bm F\cdot\operatorname{curl}\bm G.
\end{align*}
二阶恒等式为
\[
\Delta(fg)=f\Delta g+2\nabla f\cdot\nabla g+g\Delta f,
\]
\[
\nabla\times(\nabla\times\bm F)
=\nabla(\nabla\cdot\bm F)-\Delta\bm F,
\]
此处 $\Delta\bm F=(\Delta F_x,\Delta F_y,\Delta F_z)$ 为固定直角坐标分量的 Laplace 算子；曲线坐标中不能忽略基向量的变化。
\subsection{极坐标下的梯度、散度与Laplace算子}\index{极坐标下的梯度、散度与Laplace算子|hyperpage}
在二维极坐标 $x=r\cos\theta,y=r\sin\theta$ 中，取单位基
\[
\bm e_r=(\cos\theta,\sin\theta),\qquad
\bm e_\theta=(-\sin\theta,\cos\theta).
\]
对 $\bm F=F_r\bm e_r+F_\theta\bm e_\theta$，在 $r>0$ 处有
\begin{align*}
\nabla f&=f_r\bm e_r+\frac1r f_\theta\bm e_\theta,\\
\operatorname{div}\bm F
&=\frac1r\frac{\partial(rF_r)}{\partial r}
+\frac1r\frac{\partial F_\theta}{\partial\theta},\\
(\operatorname{curl}\bm F)_z
&=\frac1r\left(\frac{\partial(rF_\theta)}{\partial r}
-\frac{\partial F_r}{\partial\theta}\right).
\end{align*}
\par\pagebreak[0]\begin{samepage}
标量 Laplace 算子为
\[
\Delta f=\frac1r\frac\partial{\partial r}\left(r\frac{\partial f}{\partial r}\right)
+\frac1{r^2}\frac{\partial^2f}{\partial\theta^2}
=f_{rr}+\frac1r f_r+\frac1{r^2}f_{\theta\theta}.
\]
这里旋度的 $z$ 分量将平面场视为三维中的 $(F_x,F_y,0)$；$r=0$ 处应回到直角坐标或检查光滑延拓。
\par\end{samepage}\pagebreak[0]
\subsection{柱坐标下的梯度、散度、旋度与Laplace算子}\index{柱坐标下的梯度、散度、旋度与Laplace算子|hyperpage}
取 $x=r\cos\theta,y=r\sin\theta,z=z$，单位基为 $\bm e_r,\bm e_\theta,\bm e_z$；$\bm F=F_r\bm e_r+F_\theta\bm e_\theta+F_z\bm e_z$。在 $r>0$ 处，
\[
\nabla f=f_r\bm e_r+\frac1r f_\theta\bm e_\theta+f_z\bm e_z,
\]
\[
\operatorname{div}\bm F
=\frac1r\frac{\partial(rF_r)}{\partial r}
+\frac1r\frac{\partial F_\theta}{\partial\theta}
+\frac{\partial F_z}{\partial z}.
\]
旋度在该单位基下的三个分量为
\begin{align*}
(\operatorname{curl}\bm F)_r
&=\frac1r\frac{\partial F_z}{\partial\theta}-\frac{\partial F_\theta}{\partial z},\\
(\operatorname{curl}\bm F)_\theta
&=\frac{\partial F_r}{\partial z}-\frac{\partial F_z}{\partial r},\\
(\operatorname{curl}\bm F)_z
&=\frac1r\left(\frac{\partial(rF_\theta)}{\partial r}
-\frac{\partial F_r}{\partial\theta}\right).
\end{align*}
标量 Laplace 算子为
\[
\Delta f=f_{rr}+\frac1r f_r+\frac1{r^2}f_{\theta\theta}+f_{zz}.
\]
\subsection{球坐标下的梯度、散度、旋度与Laplace算子}\index{球坐标下的梯度、散度、旋度与Laplace算子|hyperpage}
沿用本章约定：$\varphi$ 为从正 $z$ 轴量起的极角，$\theta$ 为方位角，
\[
x=r\sin\varphi\cos\theta,\quad
y=r\sin\varphi\sin\theta,\quad z=r\cos\varphi.
\]
单位基与分量表示为
\begin{gather*}
\bm e_r=(\sin\varphi\cos\theta,\sin\varphi\sin\theta,\cos\varphi),\\
\bm e_\varphi=(\cos\varphi\cos\theta,\cos\varphi\sin\theta,-\sin\varphi),\quad
\bm e_\theta=(-\sin\theta,\cos\theta,0),\\
\bm F=F_r\bm e_r+F_\varphi\bm e_\varphi+F_\theta\bm e_\theta.
\end{gather*}
\par\pagebreak[0]\begin{samepage}
在 $r>0$、$0<\varphi<\pi$ 处，
\[
\nabla f=f_r\bm e_r+\frac1r f_\varphi\bm e_\varphi
+\frac1{r\sin\varphi}f_\theta\bm e_\theta,
\]
\[
\operatorname{div}\bm F
=\frac1{r^2}\frac{\partial(r^2F_r)}{\partial r}
+\frac1{r\sin\varphi}\frac{\partial(\sin\varphi F_\varphi)}{\partial\varphi}
+\frac1{r\sin\varphi}\frac{\partial F_\theta}{\partial\theta}.
\]
\par\end{samepage}\pagebreak[0]
旋度的单位基分量为
\begin{align*}
(\operatorname{curl}\bm F)_r
&=\frac1{r\sin\varphi}\left[
\frac{\partial(\sin\varphi F_\theta)}{\partial\varphi}
-\frac{\partial F_\varphi}{\partial\theta}\right],\\
(\operatorname{curl}\bm F)_\varphi
&=\frac1r\left[\frac1{\sin\varphi}\frac{\partial F_r}{\partial\theta}
-\frac{\partial(rF_\theta)}{\partial r}\right],\\
(\operatorname{curl}\bm F)_\theta
&=\frac1r\left[\frac{\partial(rF_\varphi)}{\partial r}
-\frac{\partial F_r}{\partial\varphi}\right].
\end{align*}
标量 Laplace 算子为
\[
\Delta f=\frac1{r^2}\frac\partial{\partial r}\left(r^2\frac{\partial f}{\partial r}\right)
+\frac1{r^2\sin\varphi}\frac\partial{\partial\varphi}
\left(\sin\varphi\frac{\partial f}{\partial\varphi}\right)
+\frac1{r^2\sin^2\varphi}\frac{\partial^2f}{\partial\theta^2}.
\]
特别地，对径向函数 $f=f(r)$，有 $\Delta f=f''(r)+2f'(r)/r$（$r>0$）。轴线与原点为坐标奇点，应以直角坐标中的光滑性判断能否延拓。
\subsection{保守场、势函数与路径无关}\index{保守场、势函数与路径无关|hyperpage}
设 $U\subset\mathbb R^n$ 为开连通区域，$\bm F\in C^1(U;\mathbb R^n)$。下列三个条件等价：存在势函数 $\phi$ 使 $\bm F=\nabla\phi$；任意两点间的分段 $C^1$ 路径积分与路径无关；每条分段 $C^1$ 闭路 $\gamma\subset U$ 均满足
\[
\oint_\gamma\bm F\cdot\mathrm d\bm r=0.
\]
此时固定 $\bm x_0\in U$，可取
\[
\phi(\bm x)=C+\int_{\bm x_0}^{\bm x}\bm F\cdot\mathrm d\bm r,
\qquad
\int_A^B\bm F\cdot\mathrm d\bm r=\phi(B)-\phi(A).
\]
在开单连通区域上，还与下列条件等价：
\[
\frac{\partial F_i}{\partial x_j}
=\frac{\partial F_j}{\partial x_i}\quad(1\le i,j\le n).
\]
在三维即 $\nabla\times\bm F=\bm0$；在二维即 $P_y=Q_x$。区域不满足单连通条件时，导数对称条件本身不足以保证路径无关。
若 $U$ 关于 $\bm x_0$ 星形且上述偏导对称条件成立，可直接构造
\[
\phi(\bm x)=C+\int_0^1
\bm F(\bm x_0+t(\bm x-\bm x_0))\cdot(\bm x-\bm x_0)\,\mathrm dt.
\]
\subsection{Green恒等式与多元分部积分}\index{Green恒等式与多元分部积分|hyperpage}
设 $\Omega\subset\mathbb R^n$ 为有界区域，边界分段光滑，$\bm n$ 为外单位法向量；以下函数在闭包的邻域内具有相应的连续导数。对 $u\in C^1$、$\bm F\in C^1$，
\[
\int_\Omega\nabla u\cdot\bm F\,\mathrm dV
=\int_{\partial\Omega}u\bm F\cdot\bm n\,\mathrm dS
-\int_\Omega u\,\operatorname{div}\bm F\,\mathrm dV.
\]
取 $\bm F=\nabla v$，对 $u,v\in C^2$ 得 Green 第一恒等式：
\[
\int_\Omega(\nabla u\cdot\nabla v+u\Delta v)\,\mathrm dV
=\int_{\partial\Omega}u\frac{\partial v}{\partial n}\,\mathrm dS,
\qquad\frac{\partial v}{\partial n}=\nabla v\cdot\bm n.
\]
交换 $u,v$ 并相减，得 Green 第二恒等式：
\[
\int_\Omega(u\Delta v-v\Delta u)\,\mathrm dV
=\int_{\partial\Omega}
\left(u\frac{\partial v}{\partial n}-v\frac{\partial u}{\partial n}\right)\,\mathrm dS.
\]
特别地，取实函数 $v=u$，得到能量恒等式
\[
\int_\Omega|\nabla u|^2\,\mathrm dV
=\int_{\partial\Omega}u\frac{\partial u}{\partial n}\,\mathrm dS
-\int_\Omega u\Delta u\,\mathrm dV.
\]
\section{无穷级数}
\subsection{常数项级数}\index{常数项级数|hyperpage}
\[
\sum_{n=1}^\infty a_n=S\iff S_N=\sum_{n=1}^Na_n\to S;
\qquad \sum a_n\text{ 收敛}\Longrightarrow a_n\to0.
\]
正项级数的比较判别与反常积分类似。比值判别设
\[
L=\lim_{n\to\infty}|a_{n+1}/a_n|;
\qquad\text{根值判别设 }L=\limsup_{n\to\infty}\sqrt[n]{|a_n|}.
\]
$L<1$ 时绝对收敛，$L>1$ 时发散，$L=1$ 时不能判定。
\[
\sum_{n=1}^\infty n^{-p}\text{ 收敛}\iff p>1.
\]
Leibniz 判别：$b_n\ge0$ 单调趋于零，则 $\sum_{n=1}^\infty(-1)^{n-1}b_n$ 收敛且 $|S-S_N|\le b_{N+1}$。绝对收敛蕴含收敛，反之不成立。
\subsection{积分判别、余项估计与Raabe判别}\index{积分判别、余项估计与Raabe判别|hyperpage}
设 $f:[1,\infty)\to[0,\infty)$ 连续且单调不增，则
\[
\sum_{n=1}^\infty f(n)\text{ 收敛}
\iff\int_1^\infty f(x)\,\mathrm dx\text{ 收敛}.
\]
收敛时，余项 $R_N=\sum_{n=N+1}^\infty f(n)$ 满足
\[
\int_{N+1}^\infty f(x)\,\mathrm dx
\le R_N\le\int_N^\infty f(x)\,\mathrm dx.
\]
Raabe 判别：若 $a_n>0$ 且下列极限存在，记
\[
R=\lim_{n\to\infty}n\left(\frac{a_n}{a_{n+1}}-1\right).
\]
$R>1$ 时 $\sum a_n$ 收敛，$R<1$ 时发散；$R=1$ 时不能判定。该判别可用于普通比值判别的极限等于 $1$ 的情形。
对实数 $p,q$，常用对数型级数的边界为
\[
\sum_{n=2}^\infty\frac1{n^p(\ln n)^q}\text{ 收敛}
\iff p>1\ \text{或}\ (p=1,\ q>1).
\]
\subsection{分部求和、Dirichlet与Abel判别}\index{分部求和、Dirichlet与Abel判别|hyperpage}
令 $A_0=0$、$A_k=\sum_{j=1}^ka_j$。分部求和（Abel 变换）给出
\[
\sum_{k=m}^na_kb_k
=A_nb_n-A_{m-1}b_m+\sum_{k=m}^{n-1}A_k(b_k-b_{k+1}).
\]
Dirichlet 判别：若 $\{A_n\}$ 有界，实数列 $b_n$ 单调趋于零，则 $\sum a_nb_n$ 收敛。
Abel 判别：若 $\sum a_n$ 收敛，实数列 $b_n$ 单调且有界，则 $\sum a_nb_n$ 收敛。
级数的 Cauchy 准则为
\[
\sum a_n\text{ 收敛}\iff
\forall\varepsilon>0,\ \exists N,\ \forall n\ge m\ge N,\quad
\left|\sum_{k=m}^na_k\right|<\varepsilon.
\]
\subsection{Bernoulli数与Euler--Maclaurin求和公式}\index{Bernoulli数与Euler-Maclaurin求和公式|hyperpage}
% Reference: https://dlmf.nist.gov/2.10
取 $B_1=-1/2$ 的约定，用生成函数定义 Bernoulli 数和多项式：
\[
\frac{t}{e^t-1}=\sum_{j=0}^\infty B_j\frac{t^j}{j!},\qquad
\frac{te^{xt}}{e^t-1}=\sum_{j=0}^\infty B_j(x)\frac{t^j}{j!}
\quad(|t|<2\pi),
\]
其中 $t=0$ 处按连续延拓解释。常用值为
\[
B_0=1,\quad B_1=-\frac12,\quad B_2=\frac16,\quad
B_4=-\frac1{30},\quad B_6=\frac1{42},\quad
B_2(x)=x^2-x+\frac16.
\]
设 $m<n$ 为整数、$p\ge1$ 为整数，$f\in C^{2p}[m,n]$，则
\[
\begin{aligned}
\sum_{k=m}^nf(k)
={}&\int_m^nf(x)\,\mathrm dx+\frac{f(m)+f(n)}2\\
&+\sum_{r=1}^p\frac{B_{2r}}{(2r)!}
\bigl[f^{(2r-1)}(n)-f^{(2r-1)}(m)\bigr]+R_p.
\end{aligned}
\]
它将有限求和表达为积分、端点修正与余项；截断后的误差由下一小节的精确余项控制。
\subsection{Euler--Maclaurin余项与误差界}\index{Euler-Maclaurin余项与误差界|hyperpage}
沿用上节条件，令 $\{x\}=x-\lfloor x\rfloor$，则精确余项为
\[
R_p=-\frac1{(2p)!}\int_m^nB_{2p}(\{x\})f^{(2p)}(x)\,\mathrm dx.
\]
因此
\[
|R_p|\le\frac{\max_{0\le u\le1}|B_{2p}(u)|}{(2p)!}
\int_m^n|f^{(2p)}(x)|\,\mathrm dx.
\]
\par\pagebreak[0]\begin{samepage}
特别地，$f\in C^2[m,n]$ 时，
\[
\sum_{k=m}^nf(k)=\int_m^nf(x)\,\mathrm dx
+\frac{f(m)+f(n)}2+\frac{f'(n)-f'(m)}{12}+R_1,
\]
\[
|R_1|\le\frac1{12}\int_m^n|f''(x)|\,\mathrm dx
\le\frac{n-m}{12}\max_{[m,n]}|f''|.
\]
\par\end{samepage}\pagebreak[0]
对无穷求和，需分别核验端点导数极限与余项积分的收敛，再令 $n\to\infty$。
\subsection{凝聚判别、绝对收敛与Cauchy乘积}\index{凝聚判别、绝对收敛与Cauchy乘积|hyperpage}
若 $a_n\ge0$ 单调不增，则 Cauchy 凝聚判别为
\[
\sum_{n=1}^\infty a_n\text{ 收敛}
\iff\sum_{k=0}^\infty2^ka_{2^k}\text{ 收敛}.
\]
绝对收敛级数可任意重排且和不变：对正整数集的任意排列 $\sigma$，
\[
\sum_{n=1}^\infty|a_n|<\infty
\Longrightarrow
\sum_{n=1}^\infty a_{\sigma(n)}=\sum_{n=1}^\infty a_n.
\]
设 $\sum_{n=0}^\infty a_n=A$、$\sum_{n=0}^\infty b_n=B$，定义 Cauchy 乘积
\[
c_n=\sum_{k=0}^na_kb_{n-k}.
\]
若两个级数均绝对收敛，则 $\sum c_n$ 绝对收敛且和为 $AB$；更一般地，若二者收敛且至少一个绝对收敛，则 $\sum c_n=AB$（Mertens 定理）。仅有两个级数各自收敛不能保证乘积级数收敛。
\subsection{无穷乘积与对数判别}\index{无穷乘积与对数判别|hyperpage}
设 $p_n=1+u_n>0$，部分乘积为 $P_N=\prod_{n=1}^Np_n$。本节将无穷乘积收敛定义为 $P_N\to P\in(0,\infty)$；趋于零的乘积另行注明。于是
\[
\prod_{n=1}^\infty(1+u_n)\text{ 收敛}
\iff\sum_{n=1}^\infty\ln(1+u_n)\text{ 收敛},
\]
\[
P=\exp\!\left(\sum_{n=1}^\infty\ln(1+u_n)\right),\qquad
P\in(0,\infty)\Longrightarrow u_n\to0.
\]
当 $t\to0$ 时，$\ln(1+t)\sim t$。特别地，$u_n\ge0$ 时，乘积收敛当且仅当 $\sum u_n$ 收敛；一般有正负项时，不能仅用等价无穷小替换断定条件收敛。
绝对收敛可用等价条件表述为
\[
\sum_{n=1}^\infty|\ln(1+u_n)|<\infty
\iff\sum_{n=1}^\infty|u_n|<\infty.
\]
此时乘积收敛于正的有限值，任意重排因子均不改变其值。若 $\sum u_n^2<\infty$，则由 $\ln(1+u_n)=u_n+O(u_n^2)$ 得到
\[
\prod_{n=1}^\infty(1+u_n)\text{ 收敛}
\iff\sum_{n=1}^\infty u_n\text{ 收敛}.
\]
\subsection{幂级数基本概念}\index{幂级数基本概念|hyperpage}
\[
\sum_{n=0}^\infty a_n(x-x_0)^n,\qquad
R=\frac1{\limsup_{n\to\infty}\sqrt[n]{|a_n|}}.
\]
约定 $1/0=\infty$、$1/\infty=0$。$|x-x_0|<R$ 绝对收敛，$|x-x_0|>R$ 发散，端点分别判断。在收敛区间内部可逐项求导与积分，收敛半径不变：
\[
\frac{\mathrm d}{\mathrm dx}\sum_{n=0}^\infty a_n(x-x_0)^n
=\sum_{n=1}^\infty na_n(x-x_0)^{n-1},\qquad
\int_{x_0}^x\sum_{n=0}^\infty a_n(t-x_0)^n\,\mathrm dt
=\sum_{n=0}^\infty\frac{a_n(x-x_0)^{n+1}}{n+1}.
\]
\subsection{幂级数的Abel端点定理}\index{幂级数的Abel端点定理|hyperpage}
若实数级数 $\sum_{n=0}^\infty a_n=S$ 收敛，则 $F(r)=\sum_{n=0}^\infty a_nr^n$ 在 $[0,1]$ 上一致收敛，且
\[
\lim_{r\uparrow1}F(r)=S.
\]
这与级数的 Abel 判别法不同，此处说明幂级数在收敛端点处的连续性。一般地，若收敛半径 $0<R<\infty$ 且 $\sum a_nR^n$ 收敛，则
\[
\lim_{x\uparrow x_0+R}\sum_{n=0}^\infty a_n(x-x_0)^n
=\sum_{n=0}^\infty a_nR^n.
\]
左端点可对 $\sum a_n(-R)^n$ 作相同处理。若两个端点都收敛，则原幂级数在整个闭收敛区间上一致收敛；端点求导仍需另行检查。
\subsection{常见幂级数展开式}\index{常见幂级数展开式|hyperpage}
\begin{align*}
e^x&=\sum_{n=0}^\infty\frac{x^n}{n!}, &&x\in\mathbb R,\\
\sin x&=\sum_{n=0}^\infty\frac{(-1)^nx^{2n+1}}{(2n+1)!}, &&x\in\mathbb R,\\
\cos x&=\sum_{n=0}^\infty\frac{(-1)^nx^{2n}}{(2n)!}, &&x\in\mathbb R,\\
\frac1{1-x}&=\sum_{n=0}^\infty x^n, &&|x|<1,\\
\ln(1+x)&=\sum_{n=1}^\infty\frac{(-1)^{n-1}x^n}{n}, &&-1<x\le1,\\
(1+x)^\alpha&=\sum_{n=0}^\infty\binom\alpha n x^n, &&|x|<1,\\
\arcsin x&=\sum_{n=0}^\infty\frac{(2n)!}{4^n(n!)^2(2n+1)}x^{2n+1}, &&|x|\le1,\\
\arccos x&=\frac\pi2-\arcsin x, &&|x|\le1,\\
\arctan x&=\sum_{n=0}^\infty\frac{(-1)^nx^{2n+1}}{2n+1}, &&|x|\le1,\\
\operatorname{arccot}x&=\frac\pi2-\arctan x, &&|x|\le1.
\end{align*}
其中 $\binom\alpha0=1$，$\binom\alpha n=\alpha(\alpha-1)\cdots(\alpha-n+1)/n!$。
以下列出三角函数在原点附近的常用展开；余切与余割是 Laurent 展开，含 $x^{-1}$ 项：
\begin{align*}
\tan x&=x+\frac{x^3}3+\frac{2x^5}{15}+\frac{17x^7}{315}+\cdots,
&&|x|<\pi/2,\\
\cot x&=\frac1x-\frac x3-\frac{x^3}{45}-\frac{2x^5}{945}-\cdots,
&&0<|x|<\pi,\\
\sec x&=1+\frac{x^2}2+\frac{5x^4}{24}+\frac{61x^6}{720}+\cdots,
&&|x|<\pi/2,\\
\csc x&=\frac1x+\frac x6+\frac{7x^3}{360}+\frac{31x^5}{15120}+\cdots,
&&0<|x|<\pi.
\end{align*}
取 $\operatorname{arcsec}x=\arccos(1/x)$，$\operatorname{arccsc}x=\arcsin(1/x)$，则对 $|x|\ge1$，
\[
\operatorname{arccsc}x=\sum_{n=0}^\infty
\frac{(2n)!}{4^n(n!)^2(2n+1)x^{2n+1}}
=\frac1x+\frac1{6x^3}+\frac3{40x^5}+\cdots,\quad
\operatorname{arcsec}x=\frac\pi2-\operatorname{arccsc}x.
\]
\section{函数列与函数项级数}
\subsection{逐点收敛、一致收敛与Cauchy准则}\index{逐点收敛、一致收敛与Cauchy准则|hyperpage}
设 $f_n,f:E\to\mathbb R$。逐点收敛允许 $N$ 依赖于 $x$；一致收敛要求 $N$ 对整个 $E$ 通用：
\[
f_n\rightrightarrows f\text{ 于 }E
\iff\forall\varepsilon>0,\ \exists N,\ \forall n\ge N,\ \forall x\in E,
\quad|f_n(x)-f(x)|<\varepsilon.
\]
记 $\|g\|_{\infty,E}=\sup_{x\in E}|g(x)|$，则
\[
f_n\rightrightarrows f\iff\|f_n-f\|_{\infty,E}\to0.
\]
一致收敛的 Cauchy 准则为
\[
\forall\varepsilon>0,\ \exists N,\ \forall m,n\ge N,\quad
\sup_{x\in E}|f_m(x)-f_n(x)|<\varepsilon.
\]
对函数项级数 $\sum_{n=1}^\infty u_n(x)$，令 $S_N=\sum_{n=1}^Nu_n$；其一致收敛即 $S_N\rightrightarrows S$，等价于
\[
\forall\varepsilon>0,\ \exists N,\ \forall n\ge m\ge N,\quad
\sup_{x\in E}\left|\sum_{k=m}^nu_k(x)\right|<\varepsilon.
\]
\subsection{一致收敛的M、Dirichlet与Abel判别}\index{一致收敛的M、Dirichlet与Abel判别|hyperpage}
Weierstrass $M$ 判别：若存在 $M_n\ge0$，使
\[
|u_n(x)|\le M_n\quad(x\in E),\qquad
\sum_{n=1}^\infty M_n<\infty,
\]
则 $\sum u_n$ 一致收敛，且 $\sum|u_n|$ 也一致收敛，并有余项估计
\[
\sup_{x\in E}\left|\sum_{n=N+1}^\infty u_n(x)\right|
\le\sum_{n=N+1}^\infty M_n.
\]
一致 Dirichlet 判别：若
\[
\sup_{x\in E}\left|\sum_{k=1}^nu_k(x)\right|\le M\quad(n\ge1),
\]
而 $v_n(x)\ge0$ 对每个 $x$ 随 $n$ 单调不增，且 $\sup_{x\in E}v_n(x)\to0$，则 $\sum u_nv_n$ 一致收敛。
\par\pagebreak[0]\begin{samepage}
一致 Abel 判别：若 $\sum u_n$ 一致收敛，实函数 $v_n(x)$ 对每个 $x$ 随 $n$ 单调，且 $\sup_{n,x}|v_n(x)|<\infty$，则 $\sum u_nv_n$ 一致收敛。
\par\end{samepage}\pagebreak[0]
\subsection{Dini定理与单调函数列}\index{Dini定理与单调函数列|hyperpage}
设 $K\subset\mathbb R^d$ 为非空紧集，$f_n,f\in C(K)$，$f_n$ 逐点收敛于 $f$。若对所有 $x\in K$ 均有
\[
f_n(x)\le f_{n+1}(x)\quad(n\ge1),
\]
或对所有 $x$ 均为相反的不等式，则 Dini 定理给出
\[
\sup_{x\in K}|f_n(x)-f(x)|\longrightarrow0.
\]
单调性针对序号 $n$，不要求每个 $f_n$ 随自变量单调；紧性、各项连续及极限连续均为此定理的条件。特别地，若连续非负函数级数在紧集上逐点收敛且其和连续，则该级数一致收敛。
\subsection{一致极限的连续性与逐项积分}\index{一致极限的连续性与逐项积分|hyperpage}
若 $f_n$ 在 $E$ 上连续且 $f_n\rightrightarrows f$，则 $f$ 连续（相对于 $E$ 的拓扑）。
若 $f_n$ 在 $[a,b]$ 上 Riemann 可积且一致收敛于 $f$，则 $f$ 可积，且
\[
\int_a^bf(x)\,\mathrm dx=\lim_{n\to\infty}\int_a^bf_n(x)\,\mathrm dx,
\]
\[
\sup_{x\in[a,b]}\left|\int_a^x(f_n(t)-f(t))\,\mathrm dt\right|
\le(b-a)\|f_n-f\|_{\infty,[a,b]}.
\]
因此，若各 $u_n$ 可积且 $\sum u_n$ 在 $[a,b]$ 上一致收敛，则
\[
\int_a^b\sum_{n=1}^\infty u_n(x)\,\mathrm dx
=\sum_{n=1}^\infty\int_a^bu_n(x)\,\mathrm dx.
\]
\par\pagebreak[0]\begin{samepage}
\subsection{函数列与函数项级数的逐项求导}\index{函数列与函数项级数的逐项求导|hyperpage}
设 $f_n\in C^1([a,b])$，在某点 $x_0\in[a,b]$ 处 $f_n(x_0)$ 收敛，且 $f_n'\rightrightarrows g$。则 $f_n$ 一致收敛于某个 $f\in C^1([a,b])$，并且
\[
f'(x)=g(x),\qquad
f(x)=\lim_{n\to\infty}f_n(x_0)+\int_{x_0}^xg(t)\,\mathrm dt.
\]
端点处导数取单侧意义。对级数，若 $u_n\in C^1([a,b])$、$\sum u_n(x_0)$ 收敛且 $\sum u_n'$ 一致收敛，则 $\sum u_n$ 一致收敛，且
\[
\frac{\mathrm d}{\mathrm dx}\sum_{n=1}^\infty u_n(x)
=\sum_{n=1}^\infty u_n'(x).
\]
仅有 $\sum u_n$ 一致收敛不足以保证逐项求导。幂级数在收敛区间的每个紧子区间上一致收敛，其逐项导数级数也如此，因而可在收敛区间内部逐项求导任意次。
\par\end{samepage}\pagebreak[0]
\subsection{Weierstrass逼近与Bernstein多项式}\index{Weierstrass逼近与Bernstein多项式|hyperpage}
对任意 $f\in C[a,b]$、$a<b$，Weierstrass 逼近定理给出
\[
\forall\varepsilon>0,\ \exists\text{多项式 }P,\qquad
\max_{x\in[a,b]}|f(x)-P(x)|<\varepsilon.
\]
在 $[0,1]$ 上，一组显式逼近多项式为
\[
B_n(f;x)=\sum_{k=0}^nf\!\left(\frac kn\right)
\binom nk x^k(1-x)^{n-k}\quad(n\ge1).
\]
记 $p_{n,k}(x)=\binom nk x^k(1-x)^{n-k}$，则 $p_{n,k}\ge0$，且
\begin{gather*}
\sum_{k=0}^np_{n,k}(x)=1,\qquad
\sum_{k=0}^n\frac kn p_{n,k}(x)=x,\\
\sum_{k=0}^n\left(\frac kn-x\right)^2p_{n,k}(x)=\frac{x(1-x)}n.
\end{gather*}
对 $f\in C[0,1]$，有 $\|B_n(f)-f\|_\infty\to0$。一般区间通过 $u=(x-a)/(b-a)$ 化到 $[0,1]$。
\subsection{Bernstein逼近的误差估计}\index{Bernstein逼近的误差估计|hyperpage}
令 $\omega_f$ 为 $[0,1]$ 上的连续模。对 $n\ge1$，有
\[
\|B_n(f)-f\|_\infty
\le\frac32\,\omega_f(n^{-1/2}).
\]
若 $|f(x)-f(y)|\le C|x-y|^\alpha$，$0<\alpha\le1$，则
\[
|B_n(f;x)-f(x)|
\le C\left(\frac{x(1-x)}n\right)^{\alpha/2},\qquad
\|B_n(f)-f\|_\infty\le\frac C{2^\alpha n^{\alpha/2}}.
\]
特别地，Lipschitz 常数为 $L$ 时，误差不超过 $L/(2\sqrt n)$。若 $f\in C^2[0,1]$，还可使用
\[
|B_n(f;x)-f(x)|\le\frac{x(1-x)}{2n}\|f''\|_\infty,
\qquad
\|B_n(f)-f\|_\infty\le\frac{\|f''\|_\infty}{8n}.
\]
这些给出特定多项式序列的误差界；其他插值与数值逼近公式见第\ref{chap:numerical-analysis}章。
\section{积分的渐近估计（进阶）}
\subsection{Laplace方法：非退化内点极小值}\index{Laplace方法：非退化内点极小值|hyperpage}
% Reference: https://dlmf.nist.gov/2.3
设 $a<b$，$f\in C^2[a,b]$、$g\in C[a,b]$，且 $f$ 在唯一内点 $x_0\in(a,b)$ 取得全局最小值，满足 $f''(x_0)>0$、$g(x_0)\ne0$。令
\[
I(t)=\int_a^be^{-tf(x)}g(x)\,\mathrm dx.
\]
当 $t\to+\infty$ 时，
\[
I(t)\sim e^{-tf(x_0)}g(x_0)
\sqrt{\frac{2\pi}{t f''(x_0)}}.
\]
若进一步 $f\in C^4[a,b]$、$g\in C^2[a,b]$，则可写成相对误差形式
\[
I(t)=e^{-tf(x_0)}g(x_0)
\sqrt{\frac{2\pi}{t f''(x_0)}}\,[1+O(t^{-1})].
\]
该形式要求唯一的非退化内点极小值；端点极小或 $f''(x_0)=0$ 时需使用相应的其他展开。
\subsection{Laplace方法：端点极小值}\index{Laplace方法：端点极小值|hyperpage}
设 $f\in C^1[a,b]$、$g\in C[a,b]$。若 $f$ 在 $a$ 处取得唯一全局最小值，且右导数 $f'(a)>0$、$g(a)\ne0$，则
\[
\int_a^be^{-tf(x)}g(x)\,\mathrm dx
\sim\frac{g(a)e^{-tf(a)}}{t f'(a)}\quad(t\to+\infty).
\]
若唯一全局最小值位于 $b$，且左导数 $f'(b)<0$、$g(b)\ne0$，则
\[
\int_a^be^{-tf(x)}g(x)\,\mathrm dx
\sim\frac{g(b)e^{-tf(b)}}{t|f'(b)|}.
\]
在相应条件下，若再有 $f\in C^2[a,b]$、$g\in C^1[a,b]$，两式均可将 $\sim$ 改写为主项乘 $[1+O(t^{-1})]$。
\subsection{Watson引理与代数型端点展开}\index{Watson引理与代数型端点展开|hyperpage}
设 $b>0$，$h\in C(0,b]$，在 $x\downarrow0$ 时满足
\[
h(x)=\sum_{k=0}^{N-1}a_kx^{\lambda+k-1}
+O(x^{\lambda+N-1}),\qquad \lambda>0,\quad N\ge1.
\]
则该区间上的积分绝对收敛，并且当 $t\to+\infty$ 时，
\[
\int_0^be^{-tx}h(x)\,\mathrm dx
=\sum_{k=0}^{N-1}\frac{a_k\Gamma(\lambda+k)}{t^{\lambda+k}}
+O(t^{-\lambda-N}).
\]
若 $h(x)\sim\sum_{k=0}^\infty a_kx^{\lambda+k-1}$ 为逐阶成立的渐近展开，则
\[
\int_0^be^{-tx}h(x)\,\mathrm dx
\sim\sum_{k=0}^\infty\frac{a_k\Gamma(\lambda+k)}{t^{\lambda+k}}.
\]
此处右端表示按有限阶截断得到的渐近展开，不要求该无穷级数收敛；把积分上限改为无穷时，还需增加对远端增长的控制条件。
\section{反演、隐函数与多维渐近补充}
级数反演和渐近方法参见\cite{dlmf}；多元微分采用对称双线性二阶导数。
\subsection{Lagrange--Bürmann级数反演}\index{Lagrange-Bürmann级数反演|hyperpage}
设 $\phi$ 在零点附近解析且 $\phi(0)\ne0$，方程 $w=z\phi(w)$ 在原点有唯一解析解 $w(0)=0$。对解析 $F$，
\[
F(w(z))=F(0)+\sum_{n\ge1}\frac{z^n}{n!}
\left.\frac{\mathrm d^{n-1}}{\mathrm dt^{n-1}}
\bigl(F'(t)\phi(t)^n\bigr)\right|_{t=0}.
\]
这是局部收敛展开，也可在特征零的形式幂级数环中理解；不能在未检查奇点时宣称对全部 $z$ 收敛。
\subsection{隐函数的二阶方向微分}\index{隐函数的二阶方向微分|hyperpage}
设 $F(x,y)=0$、$F\in C^2$、$F_y$ 可逆，隐函数为 $y=g(x)$。记 $g_h=Dg(x)h=-F_y^{-1}F_xh$，则
\[
D^2g(x)[h,k]=-F_y^{-1}\bigl(F_{xx}[h,k]
+F_{xy}[h,g_k]+F_{yx}[g_h,k]+F_{yy}[g_h,g_k]\bigr).
\]
各偏导在 $(x,g(x))$ 处取值。该式把一阶隐函数公式再微分，适用于向量值 $y$；标量分母形式不能直接套用于矩阵 $F_y$。
\subsection{多维Laplace主项}\index{多维Laplace主项|hyperpage}
有界域上，$\phi$ 有唯一内点全局极小 $x_0$，在其邻域 $C^2$，Hessian $H\succ0$，并与邻域外函数值有正间隙。若振幅 $a$ 连续、可积且 $a(x_0)\ne0$，则
\[
\int_\Omega a(x)e^{-\lambda\phi(x)}\,\mathrm dx
\sim a(x_0)e^{-\lambda\phi(x_0)}
\left(\frac{2\pi}{\lambda}\right)^{d/2}(\det H)^{-1/2}.
\]
当 $a(x_0)=0$ 时应写成相应的主项加 $o(e^{-\lambda\phi(x_0)}\lambda^{-d/2})$；退化极小和边界极小有不同尺度。
\subsection{Reynolds移动区域积分公式}\index{Reynolds移动区域积分公式|hyperpage}
光滑移动区域 $\Omega_t$ 由速度场 $v$ 运输，对足够光滑的 $f(t,x)$，
\begin{align*}
\frac{\mathrm d}{\mathrm dt}\int_{\Omega_t}f
&=\int_{\Omega_t}\partial_tf+\int_{\partial\Omega_t}f\,v\cdot n\,\mathrm dS\\
&=\int_{\Omega_t}(\partial_tf+\operatorname{div}(fv)).
\end{align*}
把 $\partial_tf+\operatorname{div}(fv)=0$ 代入即得随体质量守恒。边界项使用区域外法向，固定区域求导只是 $v\cdot n=0$ 的特例。
\clearpage
\end{document}
